\documentclass[../main.tex]{subfiles}
\begin{document}
%==========================================TIÊU ĐỀ TẬP========================================
\begin{table}[h]
  \begin{adjustwidth}{-1cm}{}
    \begin{tabular}{>{\centering\arraybackslash}p{6cm}|>{\raggedright\arraybackslash}p{12.5cm}}
      \multirow{5}{6cm}
      % Cột bên trái: ÉPISODE và số tập
      {\\[-40pt]
        \flaregothic\fontsize{20pt}{20pt}\selectfont \centering
        \color{eptype}HISTOIRE\color{black}\\[12pt]
        \dawnland\fontsize{36pt}{36pt}\selectfont
        \color{epnum}5
        \\[18pt]
        \flaregothic\fontsize{14pt}{14pt}\selectfont
        \color{epattr}BIENVENUE, \uline{\color{blue} Game} !
        % \flaregothic\fontsize{18pt}{18pt}\selectfont
        % \color{epattr}FIN DE TOME ? 
        \color{black}
      }
       &
      % Dòng thứ nhất: tên tập tiếng Nhật
      {\fotskip\fontsize{18pt}{18pt}\selectfont ゲームがありません。\ruby{間}{ま}\ruby{違}{ちが}った\ruby{寸}{すん}\ruby{法}{ぽう}。
      }            \\[6pt]
      % Dòng thứ hai: tên tập tiếng Anh
       & {\cafeteria\fontsize{18pt}{18pt}\selectfont There Is No Episode. Or Maybe There is, But in Another Universe?} \\[4pt]
      % Dòng thứ ba: tên tập tiếng Việt
       & {\vnmsans\fontsize{18pt}{18pt}\selectfont Cuộc phiêu lưu của Kuon ở thế giới "không-phải-là-Game"}                       \\[8pt]
      % Dòng thứ tư: Nhân vật xuất hiện
       & {\caviar{\fontsize{14pt}{14pt}\selectfont Nhân vật: \emp{kuon1} \emp{kaigou1} \emp{suzuki1}}}
      \\[18pt]
      % % Dòng cuối cùng: Thời gian diễn ra của tập (thường sẽ là ngày phát hành)
      % & {\continuum\fontsize{16pt}{16pt}\selectfont Ngày phát hành: 12/09/2021}\\[18pt]
    \end{tabular}
  \end{adjustwidth}
\end{table}
%===========================================NỘI DUNG==========================================
\color{black}

\pov{Hikawa Kuon}

\begin{center}
  \uline{Vũ trụ 256.\\(Dựa theo \textit{There is No Game - Jam Edition} của \textit{Draw Me A Pixel}.)}
\end{center}

Tôi mở mắt ra.

Một thứ gì đó vừa cứng vừa lạnh áp sát lên lưng tôi.

Cảm giác đầu tiên mà tôi nhận được. Không phải kiểu rắn chắc như nền xi măng, mà là cái gì đó rắn nhưng mịn, như nền đá cẩm thạch được mài nhẵn.

Tôi bật ho khẽ một tiếng khi cảm thấy lưng mình hơi tê, và trước khi kịp ngồi dậy thì *BỤP!* một vật thể nào đó rơi ngay cạnh tôi, suýt nữa làm tôi giật mình bật dậy ngay lập tức.

\dots Không, là một ai đó.

``Ái cha\dots'' giọng Kaigou-chan cất lên bên cạnh, đầy bực bội.

Tôi quay sang thì thấy cô nàng nằm nghiêng dưới đất, tay chống lên phần hông. Vẫn chưa kịp cất lên một lời nào thì\dots

*BỊCH!*

``\dots A!'' một tiếng rơi nhẹ khác vang lên, lần này là ở ngay trên người tôi.

``\textbf{Suzuki!?}'' tôi giật mình, bởi lẽ thân hình nhỏ nhắn ấy vừa hạ cánh gọn gàng trong lòng tôi.

``\vie{Ehehe\~{} Em xin lỗi\dots\ nhưng chỗ này ấm lắm,}'' cô bé cười gượng.

Kaigou thở dài, vừa chống tay ngồi dậy vừa lườm tôi như thể trách móc: ``Tay anh trai đa năng mà cũng làm nệm được à?''

Tôi nhướng mày đáp lại: ``\vie{Đừng nhìn tôi như thể tôi chủ động hứng con bé.}''

Sau vài giây điều chỉnh lại nhịp thở, cả ba chúng tôi cùng đứng dậy, phủi bụi còn vương trên quần áo. Đập vào mắt tôi là một màn đêm tuyệt đối. Không có trần, không có sàn rõ rệt, không có nguồn sáng. Chỉ có ánh sáng lờ mờ từ chính cơ thể ba chúng tôi tỏa ra, đủ để nhìn thấy nhau.

``\vie{\dots Chúng ta vẫn còn ở The Void sao?}'' cô chị gái lên tiếng, giọng cô vẫn còn khàn khàn sau cú ngã.

Tôi nhìn quanh một lượt, không có một cái gì cả. ``\vie{Không giống lắm. The Void có lớp sương tím nhè nhẹ, dù ít hay nhiều cũng có ánh sáng ma quái. Còn đây thì\dots}'' tôi nói, vươn tay ra trước mặt, không cảm nhận được một cái gì cả, ``\vie{\dots chỉ có màu đen tuyền thôi. Không có lấy một chút sương hay mảng sáng nào. Hoàn toàn trống rỗng.}''

Cô em gái gật đầu đồng tình, mặt đầy ngơ ngác: ``\vie{Thế\dots\ chúng ta đang ở đâu vậy?}''

``\vie{Anh không chắc. Nhưng để biết rõ hơn, ta thử đi một vòng xem nào.}''

Chúng tôi chia ra ba hướng để thử thăm dò không gian xung quanh.

Tôi đi về bên trái, cẩn thận từng bước, cho đến khi tay tôi đụng phải vào một bức tường vô hình. Tay tôi như chạm vào mặt kính trong suốt, hoàn hảo, không có lấy một gờ nối. Cảm tưởng như một sức mạnh vô hình nào đó ngăn cản tôi đi xa hơn về bên kia.

``\vie{\dots Có tường.}''

``\vie{Ái-cha- Đụng rồi!}'' Kaigou-chan từ bên phải kêu lên. ``\vie{Đá vào cái gì cứng như bê tông vậy. Đau ghê!}''

Tôi khẽ nhăn mặt, rút tay lại, trong khi Kaigou càu nhàu nhìn chân mình, xoa xoa bên mắt cá.

``Ơ\dots\ ái da!'' một tiếng *CỘP* nhỏ vang lên từ phía sau, và tôi quay lại thì thấy Suzuki ôm trán, mặt nhăn nhó: ``\vie{Em đi thử phía sau, nhưng lại bị đụng đầu\dots\ ở đây cũng có tường nữa rồi.}''

``\vie{\hl{Vậy là bốn phía đều có vách ngăn,}}'' tôi lẩm bẩm.

``\vie{Khác với The Void rồi đấy,}'' cô chị lớn nói. ``\vie{Chỗ đó là một không gian mở vô định, còn đây như cái phòng đen đặc vậy.}''

Tôi gật đầu. ``\vie{Một căn phòng vô hình, không tường, không cửa, nhưng có giới hạn.}''

``\vie{Em thấy căn phòng này có vẻ\dots\ được tạo ra để giam giữ gì đó,}'' Suzuki nói khẽ, hơi ngập ngừng.

Không ai nói gì thêm.

Chúng tôi thử tiếp tục đi quanh nơi tối tăm này, nhưng khoảng cách đi được chưa đầy chục bước thì đã chạm tường. Có cảm giác như không gian này nhỏ hơn tôi tưởng, như một cái buồng kín, không có điểm tựa, không có hướng đi, và cũng không có thời gian trôi qua.

Tôi ước chừng không gian này rộng đâu đó khoảng hai lần, dựng thành một hình chữ nhật, như một căn phòng có độ rộng tiêu chuẩn. Nhưng ngoài ra, không có bất cứ một cái gì cả.

``\vie{Vô vọng nhỉ,}'' Kaigou-chan thở dài. ``\vie{Tôi bắt đầu cảm thấy như mình đang bị nhốt trong hộp vậy.}''

``\vie{Giống như bị nhốt vào một buồng giam nghiên cứu cách ly khỏi xã hội,}'' tôi nói, mắt nhìn quanh. ``\vie{Tĩnh lặng và chết chóc.}''

``\vie{Hể!? C-Chết chóc!?}'' cô em gái lặp lại, có phần hoang mang.

``\vie{Không phải theo nghĩa đen,}'' tôi trấn an, ``\vie{kiểu như thể thế giới này chưa từng tồn tại sự sống.}''

Cả ba chúng tôi im lặng thêm một lúc, rồi đồng loạt thở dài đầy chán nản.

``\vie{\dots Ngồi xuống thôi,}'' Kaigou-chan nói, vỗ vỗ vào chân. ``\vie{Đi lòng vòng thế này chẳng để làm gì.}''

Tôi gật đầu. ``\vie{Đồng ý. Tốt nhất cứ chờ xem chuyện gì sẽ xảy ra tiếp theo vậy.}''

Suzuki kéo áo tôi nhẹ một cái: ``\vie{Anh nghĩ\dots\ chúng ta sẽ bị kẹt ở đây mãi sao?}

Tôi nhìn sang hai chị em, rồi lắc đầu.

``\vie{Tôi không nghĩ thế. Thứ gì đưa chúng ta tới đây\dots\ có lẽ là có mục đích nào đó. Và chúng ta\dots\ vẫn chưa kết thúc chuyến đi này đâu.}''

Ở giữa màn đen, giữa căn phòng không lối thoát với bốn bờ tường vô hình, giữa tiếng im lặng sâu thẳm nhất mà tôi từng nghe trong đời, chúng tôi ngồi xuống.

Dù gì cũng là một dịp để ba chúng tôi có thể ngơi chân thực sự sau một chuyến đi bộ rất dài.

Tôi khẽ rút điện thoại trong túi ra. Màn hình sáng lên một cách chậm chạp, như thể thiết bị cũng đang lười biếng phản hồi với tôi. Con số hiện lên\dots

8:15.

Tôi nheo mắt lại. ``\vie{\hl{Lạ thật, trước đó nó đứng ở 8:13 mà\dots}}'' tôi lẩm bẩm.

Ngay lúc ấy, Kaigou-chan cũng bật điện thoại lên, rồi cau mày. ``\vie{Của tôi là\dots\ 0:13. Hình như lúc nãy vẫn còn dừng ở 0:11.}''

Tôi liếc qua, rồi kiểm tra đồng hồ đeo tay của mình. Kim giây đang di chuyển.

Chậm rãi. Từng nhích, từng nhích một.

``\vie{Thời gian đang trôi,}'' tôi nói khẽ. ``\vie{Khác với The Void.}''

Cô nàng nhìn tôi, đôi mắt thoáng ngạc nhiên. ``\vie{Ờ nhỉ\dots\ trong The Void, cả điện thoại lẫn đồng hồ đều đứng yên như chết. Ở đây thì\dots}''

``\vie{Ít nhất có dòng thời gian. Tức là nơi này\dots\ có quy luật hơn.}''

Suzuki nghiêng đầu: ``\vie{Vậy\dots\ đây là một thế giới thật ạ?}''

``\vie{Cũng chưa chắc,}'' cô đáp, rồi nhíu mày. ``\vie{Nhưng là một tầng không gian có logic rõ ràng hơn The Void. Không có sự hỗn độn nào cả.}''

Tôi nhanh tay điều chỉnh lại điện thoại của cả ba người, đồng bộ theo giờ của đồng hồ đeo tay mình. Vừa xong, màn hình sáng lên rọi nhẹ vào gương mặt ba người chúng tôi trong không gian tối tăm.

``\vie{Tám giờ mười lăm}'' tôi nhẩm. ``\vie{Thời gian tiếp tục. Vấn đề là nó sẽ dẫn tới điều gì?}''

Kaigou-chan không đáp, chỉ lặng lẽ nhét điện thoại vào túi quần, rồi cùng tôi quay về điểm chờ ban đầu, trong khi Suzuki vẫn còn đang suy nghĩ đâu đó xa xăm.

\ct

\dots

\dots

Bầu không khí im lặng vẫn bao trùm, đặc quánh và ngột ngạt. Suzuki ngồi thụp xuống một góc, gác cằm lên đầu gối, còn tôi và Kaigou-chan chỉ biết nhìn quanh trong vô vọng, hy vọng một điều gì đó sẽ xảy ra.

Và rồi, nó thật sự xảy ra.

Từ hư không phía trước, giữa màn đen tuyền như mực, hai biểu tượng lấp lánh bất ngờ hiện lên. Cả ba người chúng tôi đều ngẩng đầu theo phản xạ.

``\vie{Một\dots\ lá cờ?} cô chị lớn lẩm bẩm.

Tôi nheo mắt. Hai lá cờ dạng pixel - một bên là cờ kết hợp giữa Mỹ và Anh, nửa nọ nửa kia; lá còn lại là cờ Pháp. Chúng lơ lửng trước mặt chúng tôi như một lựa chọn của trò chơi điện tử.

``\vie{Giống giao diện chọn ngôn ngữ ở mấy trò chơi cho máy PS1 ngày xưa quá,}'' tôi nhận xét.

Cô em gái tiến lên, nghiêng đầu: ``\vie{Chẳng lẽ\dots\ đây là một loại `menu'? Tức là nơi này giống như\dots\ trò chơi ạ?}''

Cô chị lắc đầu: ``\vie{Không chắc, nhưng nếu đúng thế thì chúng ta phải chọn một ngôn ngữ.}''

Tôi nhìn hai lá cờ, chần chừ di đi di lại giữa chúng. Tôi đang có cảm tính rằng đây là thế giới trong một trò chơi điện tử nào đó, và việc chọn lá cờ này có nghĩa là chọn ngôn ngữ để ta có thể tương tác.

Kaigou-chan bất ngờ quay sang hỏi: ``\vie{Cậu biết tiếng Pháp à?}''

``\vie{Tôi biết sơ sơ thôi,}'' tôi đáp. ``\vie{Không đủ để đọc văn bản phức tạp hay giao tiếp cao siêu gi đâu. Hai người thì sao?}''

``\vie{Em cũng không biết một chữ tiếng Pháp nào luôn,}'' Suzuki giơ tay như thú nhận lỗi. ``\vie{Em học tiếng Nhật là đủ mệt rồi\dots}''

``\vie{Vậy thì tiếng Anh nhé,}'' Kaigou gật đầu xác nhận, rồi huých nhẹ tay vào tôi. ``\vie{Cậu là đàn ông mà. Mạnh mẽ lên chứ!}''

Tôi thở dài, giơ bàn tay lên cao. ``\vie{Chỉ là chọn ngôn ngữ thôi mà\dots}''

Tôi dùng bàn tay chạm vào lá cờ nửa Mỹ nửa Anh. Ngay lập tức, biểu tượng đó chập chờn vài giây, rồi tan biến, kéo theo cả lá cờ Pháp.

Không khí vẫn im ắng. Thế rồi\dots

*Ù Ù Ù*

Một âm thanh bass bị méo mó kỳ quặc vang lên, như tiếng loa rè đang cố thốt ra điều gì. Chúng tôi giật nảy mình.

``\vie{Nghe sợ thế\dots}'' cô chị lớn lùi lại nửa bước, ôm lấy tay em gái mình theo phản xạ.

Tôi nheo mắt: ``\vie{Để xem nào\dots}''

Một dòng chữ hiện lên, được viết bằng tiếng Anh:

\begin{center}
  \uline{Not a game by\\DRAW ME A PIXEL}
\end{center}

Bên cạnh dòng chữ là một hình ảnh mô phỏng một cuốn sách đang mở, còn ở chính giữa màn hình là một khối lập phương xanh nước biển trôi lơ lửng.

\img{ting1}

``\eng{Draw Me A Pixel?}'' cả hai chị em đọc to. ``\vie{Sao tên nghe lạ thế\dots}''

``\vie{Tôi nghĩ chắc là tên nhà làm game đấy,}'' tôi đáp, rồi chống hông, khẽ gật đồng tình với họ: ``\vie{Cơ mà\dots\ tên dị vãi.}''

Tiếp theo, một dòng chữ to hơn hiện ra, từng chữ cái được đặt trong ô vuông sáng bóng màu xanh (cũng được pixel hóa):

\begin{center}
  T H E R E   I S   N O   G A M E
\end{center}

\img{ting2}

Suzuki ngẩn ra. ``\vie{Không\dots\ có trò chơi?}''

Tôi gật đầu rất khẽ, ánh mắt không rời khỏi màn hình trước mặt. Cái tiêu đề này, tôi đã từng thấy qua rồi.

``\vie{Hình như tôi đã từng chơi con game này rồi\dots}'' Kaigou-chan chợt nói, giọng trầm ngâm. ``\vie{Có liên quan đến một giọng nói, rồi point-and-click nữa\dots}''

``\vie{Em nhớ giọng của ông ta ồm ồm, rồi còn khản đặc nữa,}'' cô em gái của cô ấy reo lên.

Tôi trợn mắt. Đúng rồi. Giờ thì tôi nhớ ra rồi.

Nó chính là cái non-game từng gây sốt cách đây sáu năm, nơi mọi hành động phá luật đều là một phần của cốt truyện kỳ quặc, nơi nhân vật đôi lúc phá luôn cả bức tường thứ tư bằng cách tương tác với người chơi. Thậm chí\dots\ \emph{ông ta} còn biết người chơi đang có cái gì. Giống như thể là một thực thể riêng biệt được tạo ra không qua lập trình.

Dòng chữ trước mặt dần rõ nét, từng khối ô xanh sáng lên như dần xác lập sự thật không thể chối cãi.

Và rồi, tiếng bass kỳ quái ấy bất chợt tắt ngấm bằng một âm thanh nghe như tiếng đĩa nhạc bị dừng đột ngột.

Mọi thứ khựng lại. Không ai nói lời nào. Ba người chúng tôi nhìn chằm chằm vào dòng chữ đó trong im lặng, khi linh cảm rằng đây mới chỉ là điểm khởi đầu.

``E hèm\dots''

Một giọng nói trầm khàn, như bị nén qua lớp lọc âm méo mó, bất chợt vang lên từ hư vô. Cả ba chúng tôi đều giật mình, theo phản xạ mà nhìn quanh, nhưng chẳng có ai cả.

Người đó hắng giọng, rồi nói tiếp, giọng lộ rõ chất Pháp dày đặc: ``\eng{Xin chào người dùng! À khoan\dots\ ở đây có\dots\ hai người à?}''

Tôi và Kaigou-chan đứng hình.

``\vie{Khoan đã\dots}'' tôi lẩm bẩm \vie{Cái giọng này\dots}''

``\vie{Là cái giọng\dots}'' cô nàng tiếp lời tôi, mắt tròn xoe, ``\vie{trong cái non-game đó!}''

Suzuki ngơ ngác, nhưng ánh mắt vẫn dõi theo tôi: ``\vie{Ý hai người là nhân vật chính trong game kia ạ? Cái người không chịu thừa nhận là game đó hả?}''

``\eng{Non-game,}'' ông ta lập tức chữa lại. ``\eng{À, với cả tôi không thấy cô ở đó, xin lỗi.}''

``\vie{Ừ,} tôi khẽ gật đầu, ``\vie{chính ông ta\dots}''

Một luồng ớn lạnh chạy dọc sống lưng tôi, vừa buồn cười, vừa rợn người. Giọng nói này: trầm, méo, và cực kỳ đặc sệt kiểu Pháp\dots\ không thể nhầm lẫn được.

Tôi lên tiếng, nửa dè chừng nửa khẳng định: ``\eng{Ông đang nói chuyện với chúng tôi à?}''

``\eng{Ờm\dots\ đúng,}'' giọng đó đáp lại. ``\eng{Tôi không nhận ra là có ba người ở đây.}''

Kaigou-chan nhướng mày, khoanh tay: ``\eng{Ông đang chọc tôi à? Mắt ông mù chắc?}''

``\eng{Ờ thì\dots\ tôi không có mắt,}'' giọng đó phản bác, ``\eng{nhưng mà tôi dùng cảm biến để phát hiện- mà này!}''

Cô chị lớn phá lên cười khúc khích khi bị cô ấy lạc qua chủ đề khác. Tôi chỉ biết thở dài bất lực, còn Suzuki thì che miệng cười nhẹ, ánh mắt vẫn lấp lánh vì tò mò.

Cô nàng dịu giọng lại: ``\eng{Xin lỗi, xin lỗi. Vậy ông là ai vậy?}''

``\eng{Các người có thể gọi tôi là\dots\ `Game'.}''

``\eng{Game\dots?}'' cô bé nghiêng đầu. ``\eng{Nhưng mà ông nói đây không phải trò chơi mà?}''

Ông ta thở dài, rồi đáp lại với giọng chán nản, như thể đây không phải lần đầu gặp phải tình huống nào: ``\eng{Tôi không thể nghĩ tên nào khác tốt hơn. Xin lỗi.}''

``\eng{Vậy\dots\ ông ở đâu?}'' Kaigou-chan hỏi. ``\eng{Sao tôi chỉ nghe thấy giọng của ông mà chẳng thấy bóng dáng đâu hết?}''

``\eng{Tôi không có hình dạng vật lý như con người,}'' Game đáp. ``\eng{Vốn dĩ tôi chỉ là một chương trình thôi.}''

``\eng{Kiểu như\dots\ A.I?}'' tôi hỏi lại.

``\eng{Có thể gọi như vậy.}''

Cả ba chúng tôi đều khẽ gật gù. ``\eng{Ra thế\dots}''

Tôi hỏi tiếp, đi thẳng vấn đề: ``\eng{Vậy chúng tôi đang ở đâu?}''

Giọng kia đáp ngay, như thể đó là điều hiển nhiên: ``\eng{Tất nhiên là các cậu đang ở trong trò chơi rồi.}''

Kaigou-chan chau mày, ngồi chống tay lên hông, lập tức vặn lại: ``\eng{'Ý ông là  không-phải-trò-chơi, đúng không?}''

Game giật mình ngập ngừng nói nhanh: ``\eng{À, đúng, đúng! Không-phải-trò-chơi! Cảm ơn vì đã nhắc!}''

Tôi toát cả mồ hôi trán. Suzuki thì rụt cổ lại, thì thầm: ``\vie{\hl{Ông ta lơ mơ thật luôn à\dots}}

Ông ta hắng giọng một lần nữa, rồi quay trở lại với giọng nửa cảm thông (dù tôi biết là giả), nửa như muốn làm cho xong chuyện: ``\eng{Nhắc mới nhớ, tôi có tin xấu cho các cậu đây.}''

Tôi thở dài, nói luôn: ``\eng{Để tôi đoán. Ông định nói là `ở đây chẳng có game nào cả' đúng không?}''

``\eng{\dots Thực ra là, ở đây chẳng có- Ồ, wow. Cậu đọc được suy nghĩ của tôi đấy.}''

Cô chị lớn lườm nhẹ: ``\eng{Cái dòng chữ to đùng treo lơ lửng thế kia là bọn này biết rồi.}''

``\eng{Ừ phải. Dù sao thì, tôi hy vọng các cậu không quá thất vọng.}''

Cô nàng bĩu môi, giọng đậm tính mỉa mai: ``\eng{Ừa, phải.}'' Tôi thì nhún vai: ``\eng{Đoán được trước rồi.}''

``\eng{Có vẻ ở đây không có gì giúp ích nhỉ,}'' Suzuki nói, giọng bắt đầu chùng xuống.

``\eng{Tôi e là vậy,}'' Game đáp. ``\eng{Các cậu có thể xem TV, ra ngoài chơi, đọc sách\dots}''

``\eng{Làm sao mà chúng tôi làm được mấy chuyện đó khi bị kẹt ở đây?}'' cô bé hỏi lại với giọng thất vọng. ``\eng{Ở đây cũng chẳng có gì để làm nữa.}''

``\eng{À, ừ, phải, tôi quên mất. Tôi cũng đang định nói là đòi hoàn tiền, nhưng mà\dots}''

``\eng{Tôi nhớ là cái tựa không-phải-là-game này miễn phí mà?}'' Kaigou-chan nheo mắt.

``\eng{\dots Ừ\dots}'' Game ậm ừ, rồi quay ra lẩm bẩm: ``\eng{\hl{Làm sao mấy người có thể đoán được tôi nói gì\dots?}}''

Suzuki giờ mới nhỏ giọng hỏi: ``\eng{Vậy là\dots\ không có cách nào rời khỏi nơi này sao, Game-san?}''

Game đáp, nghe có chút tiếc nuối: ``\eng{Nếu các cậu đang định tìm cách thoát khỏi đây, thì\dots\ tôi e là không khả thi đâu. Ở đây chỉ có sự chán nản và vô định thôi.}''

Tôi thở dài, mắt hướng về màn đêm xung quanh:
``\vie{Vậy là chúng ta sẽ bị kẹt ở đây\dots}''

Kaigou-chan cũng thả người ngồi xuống bên cạnh tôi, tay chống cằm, gật gù: ``\vie{Dễ là như vậy nhỉ\~{}}''

Suzuki ngồi xuống phía đối diện, nhìn cả hai người chúng tôi, vẻ mặt tuy còn lo lắng nhưng vẫn giữ được chút hy vọng mong manh: ``\vie{Onii-sama, Onee-sama\dots\ chúng ta sẽ bị kẹt ở đây chứ?}''

Tôi quay sang cô bé, nở một nụ cười mệt mỏi, nhưng vững vàng: ``\vie{Dĩ nhiên là không. Anh chị có cách.}''

Ngay khi tôi dứt câu, Game nói, có vẻ như đang cố chấm dứt cuộc trò chuyện: ``\eng{Vậy, cứ để mặc tôi đi nha. Tôi sẽ rất cảm kích đấy. Dù gì cũng cảm ơn ba người và\dots\ tạm biệt.}''

*CẠCH*

Một tiếng khô khốc như ai đó tắt míc-rô. Không gian lại chìm vào yên lặng.

Tôi thở ra, tay chống cằm nhìn vào khoảng không trước mặt như thể đang nhìn một màn hình vô hình nào đó. ``\vie{Có vẻ như chúng ta đang thật sự ở trong thế giới của There Is No Game,}'' tôi lẩm bẩm.

Cô nàng quay đầu sang nhìn tôi: ``\vie{Dường như là vậy rồi.}'' Suzuki cũng đồng ý với tôi.

``Mà thật ra, ả hai người đều biết trò này rồi, đúng không?'' tôi liếc sang hai chị em, hỏi một câu thật lòng.

Cô em gái khẽ gật đầu, nụ cười nhẹ hiện lên: ``\vie{Vâng. Em từng thấy Onee-sama chơi trò này hồi chị ấy lên lớp Chín trên laptop mượn hàng xóm ạ. Em ngồi cạnh xem mà cứ tưởng là bị lỗi thật, ai dè nó hay ghê.}''

``\vie{Xem chừng cả ba chúng ta đều biết cách thứ này vận hành rồi nhỉ?}''

Tôi khoanh tay, trầm ngâm: ``\vie{Chắc vậy\dots\ Nhưng tôi vẫn muốn thử một điều.}''

``\vie{Điều gì?}''

``\vie{Kiểm tra xem ông ta có nói đúng y như bản gốc không. Nếu đúng, chúng ta có thể dự đoán được hành vi của ông ta, từ đó điều hướng cốt truyện.}''

Kaigou-chan nghiêng đầu: ``\vie{Giống như chơi lại một game đã từng chơi, nhưng lần này là trong game thật?}''

``\vie{Không phải \emph{game},}'' tôi chỉnh lại, mỉm cười, \vie{\emph{non-game} mới đúng.}''

Suzuki phá lên cười nhỏ: ``\vie{Onii-sama nói chuẩn luôn đó.}''

``\vie{Thế thì ta làm gì bây giờ?}'' cô chị lớn duỗi chân, ngả lưng ra phía sau. ``\vie{Ngồi đợi ông ta mở mồm hả?}''

``\vie{Còn cách nào nữa đâu,}'' tôi đáp.

Chúng tôi im lặng vài giây. Bầu không khí tuy không ngột ngạt như The Void, nhưng lại phảng phất một sự trống rỗng quen thuộc. Kaigou-chan là người phá tan sự yên tĩnh đầu tiên: ``\vie{Cậu có chuyện gì vui vui kể không? Đỡ chán chút đi.}''

``\vie{Có cả rổ. Muốn nghe chuyện gì?}'' tôi ngước mắt lên trên, rồi nhìn lại về phía hai chị em.

``\vie{Chuyện gì cũng được. Nhưng phải thật hài vào nha.}''

Cô em gái góp lời, vẫy tay: ``\vie{Vậy thì anh thử kể chuyện làm anh muối mặt nhất đi, Onii-sama!}''

Tôi cười chát: ``\vie{Bao nhiêu chủ đề không chọn lại chọn cái đó. Thôi được rồi.}''

Tôi hít một hơi thật sâu, và ngay lúc tôi chuẩn bị kể, một giọng nói quen thuộc lại vang lên, lần này đầy vẻ luống cuống: ``\eng{À\dots\ quên nói với ba người. Đừng có chạm vào tiêu đề đấy. Nó vẫn chưa khô hẳn đâu.}''

Chúng tôi đồng loạt quay đầu nhìn về phía tiêu đề ``THERE IS NO GAME'' đang lơ lửng cách mặt đất một đoạn. Kaigou-chan phồng má, bĩu môi đầy mỉa mai: ``\eng{Mồ\~{}! Tôi tưởng là ông muốn chúng tôi kệ ông cơ mà?}''

``\eng{Tôi xin lỗi,}'' Game đáp với giọng y như đang chống chế, ``\eng{chỉ là tôi muốn nhắc nhở ba người thôi.}''

Tôi nhún vai: ``\eng{Cũng chẳng có gì để làm mà. Thử đi loanh quanh thì đập đầu vào tường vô hình. Ngồi kể chuyện thì bị ngắt lời.}''

Ông ta thở dài lấy một tiếng, rồi giọng của Game hạ xuống như thì thầm, lẩm bẩm nghe mà gai người: ``\eng{Đừng động vào tiêu đề đấy. Đặc biệt là chữ \emph{O}.}''

Cả ba chúng tôi nghiêng đầu.

``\eng{Chữ O?}'' cô nàng lặp lại.

Tôi chống cằm suy nghĩ: ``\vie{Tôi nghĩ là ông ta đang ám chỉ\dots\ đúng cái trò chúng ta từng chơi. Trong game đó, cũng có đoạn bắt đầu tương tự, rồi phải làm gì đó với chữ O\dots}''

Suzuki reo lên, như vừa nhớ ra điều gì đó quan trọng: ``\vie{À, em nhớ rồi! Cái chữ O sẽ rơi xuống nếu mình cứ động vào liên tục!}''

Tôi ngước nhìn chữ O đang lơ lửng giữa không trung, tách biệt hẳn ra khỏi các chữ còn lại của dòng tiêu đề. Nhưng khổ nỗi, ở đây chẳng có gì để mà với tới hay tương tác cả.

Tôi nhìn xuống đôi dép quai hậu đang mang và lẩm bẩm: ``Thử xem vậy\dots''

Tôi tháo một chiếc dép, cân nhắc vài giây như đang ngắm một trái bóng ném, rồi ném thẳng nó lên chữ O với một cú phóng đầy chính xác.

*PẶC!*

Một tiếng động nhẹ phát ra như kim loại chạm vào kính. Một vài hạt bụi mờ rơi xuống, và chữ O tụt xuống khoảng một phân.

Kaigou-chan nhướng mày ngạc nhiên: ``\vie{Ể? Có tác dụng thật kìa?}''

Không nói không rằng, cô cũng tháo chiếc dép của mình ra, xoay xoay một vòng rồi ném thẳng lên.

*BỐP!*

Lần này, chữ O rung lắc mạnh hơn, tụt xuống rõ rệt gần ba phân. Tôi nhìn cô ấy, nửa tò mò nửa ngán ngẩm: ``\vie{Cô ném thế không sợ hỏng dép à?}''

Cô ấy hất tóc, nhếch môi đầy tự tin. Nhưng trước khi cô kịp trả lời, Suzuki chen ngang, hai tay chống nạnh: ``\vie{Cái dép đó là cái bền nhất mà mẹ chúng em để lại đó! Chọi đá còn không sứt, huống chi là chữ cái\~{}!}''

``\vie{Vậy à,} tôi bật cười. ``\vie{Em có muốn thử không, Suzuki?}''

``\vie{Có ạ\~{}!}''

Cô bé chẳng nói chẳng rằng, giật luôn chiếc giày vải bên trái, nhắm thẳng vào chữ O.

*COONG!*

Lại thêm một cú trúng đích. Chữ Otụt thêm vài phân, giờ đã lơ lửng thấp hơn đáng kể.

``\vie{Thích quá\~{}}'' Suzuki cười toe, tháo luôn chiếc giày còn lại. ``\vie{Em chưa từng thấy ai cho phép mình ném chữ cái bao giờ!}''

Kaigou-chan bật cười, hét lớn: ``\textbf{\vie{Đại hội tạt chữ bắt đầu!!}}''

Không khí đột nhiên vui vẻ kỳ lạ như mùa hội thao. Giữa một không gian tối mịt mù, ba người chúng tôi đang thi nhau ném giày dép như thể đang chơi tạt lon dưới sân trường cấp hai. Những cú ném tiếp nối nhau như pháo hoa: *BỐP! *KENG!* CHÁT!* *BỤP!* và chữ O mỗi lúc lại tụt xuống thêm một chút.

Tới lần thứ ba mà cô chị lớn ném trúng, lần này là một cú quăng đầy lực sau khi xoay tay như đang ném đĩa bay, chữ O cuối cùng rơi rụng khỏi dòng tiêu đề.

*ỤP!*

Một tiếng trầm đục vang lên như tiếng bass bị méo. Chúng tôi ngẩng lên nhìn, thấy chữ O giờ nằm dưới đất, khẽ rung lắc như thể đang ngượng ngùng vì bị vạch trần.

\img{ting3}

Bất chợt, giọng nói trầm mặc quen thuộc thảng thốt vang lên: ``\eng{Mấy người vừa làm gì vậy!? Phá hỏng hết cả tiêu đề của tôi rồi!}''

Cô chị lớn chống hông, nhướng mày đầy khiêu khích: ``\eng{Ai bảo ông nói ra làm chi. Ông tự khai mà.}''

Tôi lắc đầu, thở dài: \vie{Dù thế thì cũng không có nghĩa là nên chọc tức người ta đâu, Kaigou-chan à.}''

Suzuki che miệng cười khúc khích: ``\eng{Onee-sama xấu tính thật á!}''

Game rền rĩ: ``\eng{Ặc. Tại cái miệng mấp máy của tôi đây mà! Dù sao thì, đưa cái chữ đó quay trở về chỗ cũ ngay!}''

Tôi ngước nhìn trần, thấy khoảng trống nơi chữ O vừa rơi xuống, rồi lại nhìn cái chữ đang nằm bẹp dưới sàn. ``\eng{Nó ở tít trên kia. Làm cách nào mà ông bảo đưa lại?}''

Người chị lớn chép miệng: ``\eng{Đúng đó. Ở đây còn chẳng có cái gì để trèo lên.}''

``\eng{Hay lấy giày xếp chồng lên nhau?}'' cô em gái đề xuất, dù trong lòng cũng biết nó khá vô lý.

``\eng{Tôi không cần biết!}'' giọng ông ta cáu lên. ``\eng{\textbf{Làm đi! Ngay và luôn!}}''

Tôi thở dài một tiếng, cúi xuống nhặt chữ O đang nằm trơ trọi dưới đất. Mặt chữ trơn nhẵn, nhưng chẳng có tí dính nào ở mặt sau. Làm thế nào mà ông ta lắp được tiêu đề trên đó được nhỉ?

Tôi thử đưa tay đặt nó trở lại vị trí trống trên dòng chữ tiêu đề đó. Ngẩng đầu, nhắm chừng chiều cao, tôi rướn người một chút, tay đưa cao hết cỡ và\dots\ dính được vào đúng chỗ đó.

``\vie{Có vẻ như được rồi đấy.}''

Tôi nhẹ thở ra, từ từ rút tay lại.

*BỤP*

Chữ O rơi lại xuống đất như thể thách thức tôi.

Kaigou-chan khoanh tay, thở dài, mắt híp lại như đang trêu tức tôi: ``\vie{Cậu vụng về thật đấy.}''

Tôi gãi đầu, nhìn sang cô: ``\vie{Ờ thì\dots\ thật ra tôi cũng công nhận. Nhưng mà hình như cái này hết keo rồi hay sao ấy.}''

``\vie{Hoặc là do cậu thôi.}''

Suzuki nhìn tôi đầy cảm thông, quay sang chị mình: ``\vie{Vậy giờ ta phải làm gì đây, Onee-sama?}''

Cô chị đứng khoanh tay một hồi lâu, rồi đập tay như vừa nghĩ ra một sáng kiến: ``\vie{Suzu, em thử đi. Cậu ấy sẽ cõng em lên. Em khéo tay hơn Kuon-kun mà.}''

``Ơ, này\dots!'' tôi chớp mắt, chưa kịp phản ứng thì cô em gái đã vui vẻ nhảy lên lưng tôi như thể đây là trò chơi cưỡi ngựa. Tôi hơi loạng choạng nhưng vẫn giữ vững được. ``\vie{Đừng có bất ngờ nhảy lên anh chứ\dots}''

``\vie{Rồi, em thử nha!}'' Suzuki rướn người, cẩn thận áp chữ O lên đúng chỗ vừa rơi xuống.

Cô bé giữ nó trong vài giây. Tôi cố đứng yên, tập trung như đang thi giữ thăng bằng trong chương trình vận động nào đó.

Suzuki từ từ rút tay ra, rồi\dots

*BỤP*

Chữ O lại rơi cái bịch xuống đất, lần này suýt trúng mũi tôi.

``Eeek!'' cô em trên lưng kêu lên, mím môi nhìn chữ cái đầy giận dỗi. ``\vie{Mất công em làm nhẹ nhàng thế cơ mà\dots}''

Tôi thở dài, nhặt lại chữ cái. ``\vie{Thấy chưa, tôi bảo mà. Không phải tại tôi đâu. Cái chữ này có khi còn chẳng dính với cái tiêu đề nữa.}''

Kaigou-chan khoanh tay lại, nheo mắt: ``\vie{Hay là tôi thử nhỉ?}''

Tôi ngước nhìn cô ấy, hỏi vu vơ nửa đùa nửa thật: ``\vie{Cô có muốn tôi cõng lên không? Nãy giờ cô cứ trù ẻo tôi nãy giờ.}''

Cô ấy hơi giật mình, tai đỏ lên thấy rõ: ``\vie{Ơ\dots\ khỏi đi. Có khi vai cậu không chịu nổi đâu. Tôi nặng hơn Suzu đấy.}''

\dots Tôi không nghe ra cái giọng có gì đùa cợt. Nó có một cái gì đó hơi khó nói và vì thế mới cảm giác thật. Tôi cũng chẳng ép thêm.

\pov{Hikawa Kaigou}

\il

Tôi nói thế vì thật ra\dots\ tôi không muốn ngồi lên cậu ta đâu. Không phải vì nặng hay nhẹ gì hết. Với cân nặng hiện tại của mình, Kuon-kun dư sức có thể cõng tôi lên.

Nhưng vấn đề không phải chỗ đó.

Tôi vẫn còn đang ôm nguyên cái bể nước ở bụng dưới từ nãy giờ! Vừa rồi còn đập chân vào góc tường, lại còn xoay vòng ném dép nữa chứ! Cứ mỗi lần ngồi xổm hay rướn người là áp lực nó lại dồn xuống. Nếu mà ngồi lên lưng cậu ta nữa thì đúng là đại họa.

Cái mặt tôi chắc giờ méo như bánh bao bị bóp dẹp nếu tôi ngồi lên cậu ta. Không cẩn thận thì sẽ có một ``con suối vàng'' chảy xuống từ lưng cậu ta cũng không chừng.

Không được. Tôi phải làm cách nào đó để không nghĩ tới chuyện đó nữa.

\il

``\vie{Tôi tự làm được,}'' Tôi nói, rồi giật lấy chữ O từ tay Kuon-kun, tự rướn người, nhảy nhẹ một cái để đặt chữ lại, vừa đủ để không gây áp lực lên bàng quang của tôi.

Tôi từ từ căn chỉnh nó vào chỗ trống, giữ thật chắc.

Ba giây\dots\

Bốn giây\dots\

Tôi từ từ rút tay xuống\dots

*BỘP!*

Chữ O rơi lại, lần này trúng trán tôi.

``Ái-cha\dots'' tôi ôm lấy trán tôi, khẽ rít lên vì đau.

Giọng Game lại vang lên, lần này như đầy bất lực: ``\eng{Thôi kệ đi. Tôi sẽ tự sửa lại lần tới khi có ai đó vào thế giới này vậy. Trông cách cô làm như thể đang làm xiếc ấy.}''

Tôi đông cứng lại. Rồi hét lớn, mặt đỏ như cà chua chín vì vừa tức vừa ngượng: ``\eng{\textbf{Thế sao ông không nói ngay từ đầu!?}}''

Tôi đá một phát vào chữ O đang lăn lóc dưới đất, suýt nữa đụng phải Suzu. Kuon-kun thì đang bịt miệng nhịn cười, còn cái giọng Game kia thì lại nhỏ nhẹ như thể chẳng có chuyện gì: ``\eng{Xin lỗi. Tôi quên mất.}''

Tôi quay mặt đi, vừa lẩm bẩm vừa rên rỉ hai tay quơ lấy phần hạ bộ: ``\vie{Mịa, quê quá\dots\ Không chỉ bị đem ra làm trò cười\dots\ mà còn sắp vỡ đê đến nơi rồi\dots\ Ư\~{} Buồn tè quá\~{}\dots}''

\ct

\pov{Hikawa Kuon}

Chúng tôi vẫn đứng đó, bối rối nhìn dòng tiêu đề vừa bị phá nát mà chẳng làm được gì. Không gian lại chìm vào im lặng một cách khó chịu. Từ nãy tới giờ, cái giọng  kia cứ như chơi trò mèo vờn chuột, còn chúng tôi thì chẳng khác nào ba con mèo con đang tìm đường ra khỏi một cái hộp.

Suzuki bất chợt nheo mắt, ngước nhìn lên trần nhà xám xịt kia, rồi chỉ tay: ``\vie{Ơ, nhìn kìa! Trên kia có cái gì nhô ra kìa!}''

Tôi cũng ngẩng đầu lên nhìn theo hướng ngón tay của em ấy. Thoạt đầu tôi chẳng thấy gì cả ngoài một mảng tường xám nhẵn bóng. Nhưng rồi, khi điều chỉnh lại ánh mắt và chú ý kỹ hơn, tôi thấy một thứ gì đó hơi lồi ra, trông như một khối lập phương trắng bị ép dính vào trần.

\img{ting4}

``\vie{Lúc đầu đâu có cái đó đâu nhỉ?}''

Kaigou-chan cũng nhìn theo, hơi nghiêng đầu: ``\vie{Thật luôn. Nó ở đó từ khi nào vậy?}''

Cả ba chúng tôi lần lượt thử nhón chân, nhảy nhẹ, với tay lên như ba đứa trẻ cố lấy lọ bánh quy trên kệ cao. Nhưng dù cố thế nào, cái khối trắng ấy vẫn cao hơn tầm với của bất kỳ ai trong nhóm tôi.

Cô nàng buông tay, thở dài ngao ngán: ``\vie{Chán thật. Chỗ này làm như thiết kế cho người cao hai mét không bằng.}''

Cô nàng có vẻ vẫn còn ấm ức chuyện vừa bị Game cà khịa, nên bất chợt đá mạnh cái chữ O vừa nãy đang nằm dưới đất.

*BỘP!*

``Ái da\dots'' cô ấy điếng người, khuỵu xuống ôm lấy bàn chân vừa đá của mình.

Chữ O bay là là qua không trung rồi rơi xuống, đập một tiếng rõ to xuống mặt sàn. Nhưng đúng khoảnh khắc đó, chúng tôi đều ngước mắt nhìn lên.

Khối trắng kia tụt xuống thêm một đoạn.

Tôi chớp mắt, rồi nhanh chóng nhìn lại chiếc chữ O.

``\vie{Khoan đã\dots}'' tôi lên tiếng, đầu óc bỗng sáng ra, ``\vie{hồi nãy khi chữ O rơi xuống, cái khối trắng kia nhô xuống thêm chút đúng không?}''

Kaigou-chan ngớ người: ``\vie{Ý cậu là tôi vô tình\dots?}''

``\vie{Không hẳn\dots\ Tôi nghĩ là cả ba chúng ta làm đấy. Lúc mà ta cố tra lại chữ O lên dòng tiêu đề.}''

Tôi nhặt lại chữ O, nhẹ nhàng nâng nó lên ngang tầm ngực, rồi thả xuống, đồng thời mắt ngước về khối trắng kia.

*BỤP*

Nó lại tụt xuống nữa.

Kaigou tròn mắt: ``\vie{Cái kiểu đó cũng được à!?}''

Tôi gật đầu xác nhận: ``\vie{Có vẻ như mỗi lần cái chữ này đập xuống đất thì cái kia tụt xuống thêm đấy. Cũng giống như trong cái mà ta chơi trong non-game kia thôi.}''

Suzuki nghe thế thì hào hứng hẳn lên. ``\vie{Vậy để em làm cho nhanh!}''

Không đợi ai đồng ý, cô bé cúi xuống, nâng chữ O rồi thả xuống, rồi lặp lại liên tục như đang tập tâng bóng rổ.

*BỤP*

*BỤP*

*BỤP*

Mỗi cú thả là một lần cái khối kia tụt xuống thêm một chút. Sau khoảng hơn chục cú, một tiếng gì đó vang lên như vừa bị đứt.

Khối trắng rụng hẳn khỏi trần, rơi xuống, vỡ ra thành nhiều mảnh vuông nhỏ như bọt xốp.

Ở giữa đống vỡ đó, thứ còn lại là một khối kim loại hình nón, bề mặt phủ một lớp lưới nhỏ, rõ ràng đó là một cái loa. Một chiếc loa đúng nghĩa đang phát ra sóng âm có thể nhìn thấy.

\img{ting5}

Ngay lập tức, giọng Game cất lên, lần này đầy vẻ rệu rã: ``\eng{\dots Nghiêm túc đấy, mấy người đang định phá tan bành mọi thứ à?}''

Cả ba chúng tôi đứng đó, nhìn cái loa với sự kinh ngạc.

Ông ta ngập ngừng. Hình như ông ta cũng vừa mới nhận ra cái gì đó. Trong chất giọng tưởng như thản nhiên kia, tôi lại nghe ra một chút gì đó bối rối: ``\eng{Với cả, tôi nghĩ\dots\ dùng cái biểu tượng kia cũng chẳng giải quyết được gì đâu. Thật đấy. Tốt hơn hết là cứ\dots\ để mặc nó đi.}''

Tôi và hai chị em nhìn nhau, rõ ràng trong ánh mắt cả ba là một thứ đồng thuận ngầm: ông ta đang cố đánh lạc hướng chúng tôi. Giọng đó quá giả tạo để có thể tin đợc.

Tôi nheo mắt, nhìn cái loa kia. Một phần tôi muốn tin lời ông ta, phần còn lại thì mách bảo điều ngược lại.

``\eng{Cái đó chỉ là-}''

Tôi chẳng biết đầu óc đang nghĩ gì lúc ấy, chỉ là tay tôi vô thức vươn ra, chạm nhẹ vào cái loa.

Ngay lập tức, biểu tượng sóng âm trên mặt loa lập tức chuyển thành hình chữ X, như thể tôi vừa bấm nút tắt tiếng trên điều khiển. Cùng lúc đó, Game rên lên một tiếng đầy quái dị, cứ như thể có ai đó vừa nhét khăn vào miệng ông ta giữa lúc ông ta đang luyên thuyên.

\img{ting6}

``Mmmmf-!? Mmm-mmf!''

Tôi vội rụt tay lại. Loa lại bật lên, tiếng của ông ta vang lên tức tối: ``\eng{Cậu bị điên à!? T-Tôi không thể nói được đây này!}''

Kaigou-chan khoanh tay, liếc nhìn cái loa: ``\eng{Kệ ông thôi. Ông phiền vãi nồi ra.}''

``\eng{Tôi là một giọng nói đấy!}'' Game gào lên như vừa bị xúc phạm danh dự, ``\eng{Mấy người không biết một giọng nói mà không nói được thì khổ sở thế nào đâu!}''

Cô nàng nhướng mày, nghiêng đầu hỏi hờ hững: ``\eng{Nó như thế nào? Cái chết à?}'' Một cái rên rỉ khản đặc đầy tức tối từ ông ta vang lên ngay sau đó.

Tôi thở ra, nửa bất lực, nửa bất ngờ vì cái kiểu nói chuyện trêu người của cô nàng: ``\vie{Đừng chọc ghẹo người ta nữa, cái cô này.}''

``\eng{Dù gì thì\dots}'' Game hậm hực, ``\eng{đừng có chạm vào cái nút ấy nữa!}''

Kaigou-chan nhếch môi thành một nụ cười tinh quái: ``\eng{Tôi vẫn chạm vào đấy, sao nào?}''

Nói là làm. Trước khi ông ta kịp mở miệng tiếp, cô nàng lại chạm vào cái loa lần nữa.

Lập tức, Game gào lên một tràng dài như lợn bị chọc tiết, một chuỗi âm thanh méo mó phát ra từ loa, nghe mà không nhịn nổi cười.

Tôi bật cười khúc khích, còn Suzuki thì cười phá lên đến mức suýt ngồi bệt xuống sàn. ``\vie{Tiếng ông ấy phát ra nghe hài quá!}''

``\vie{Cô còn bạo hơn tôi tưởng đó,}'' tôi lắc đầu, cười nửa ngao ngán, nửa khâm phục. Cô chị lớn chớp mắt, nở nụ cười nửa nghiêm, nửa đùa: ``\vie{Tôi mà không thế thì làm sao tồn tại được ở xã hội cũ!}''

Cái loa lại bật lên, Game giờ thì gần như gào thét: ``\eng{\textbf{Cô đang định giết tôi đấy à!?}}''

``\eng{Giết gì đâu,}'' cô nàng hất tóc đuôi ngựa ra sau, đáp tỉnh bơ. ``\eng{Tôi chỉ muốn chơi đùa với ông thôi.}''

``\eng{Chơi với chúng tôi đi, Game-san!}'' Suzuki reo lên thích thú.

``\eng{Chơi với tôi chứ gì!? Đã vậy thì\dots}''

Đột nhiên, cái loa biến mất không một dấu vết ngay trước mắt tôi.

``Ể?'' Tôi lùi lại một bước.

``\vie{\textbf{Kuon-kun! Bắt lấy cái loa đó!}}'' Kaigou-chan hét lên.

Một tiếng *BÚP!* nhẹ phát ra sau lưng tôi, cái loa bất ngờ xuất hiện ở góc phòng.

Tôi nhào tới, nhưng lại biến mất, làm tôi bắt trượt.

*BÚP!* Lần này nó lại xuất hiện bên cạnh Suzuki, nhưng vừa lúc em ấy quay sang thì nó đã biến mất.

Một trò rượt đuổi bất ngờ bắt đầu, giữa ba chúng tôi và một cái loa biết dịch chuyển tức thời. Cứ mỗi lần chạm gần được là nó lại biến mất ở chỗ khác. Đôi lúc nó còn xuất hiện quá tầm với của cả ba bọn tôi, đôi lúc thì nó ở ngay trước mắt nhưng cũng không bắt được.

\begin{figure}[H]
  \centering
  \begin{subfigure}[t]{0.45\textwidth}
    \centering
    \includegraphics[width=8cm]{../../../../Image/BookV/ting7a.png}
  \end{subfigure}
  \quad
  \begin{subfigure}[t]{0.45\textwidth}
    \centering
    \includegraphics[width=8cm]{../../../../Image/BookV/ting7b.png}
  \end{subfigure}
  \quad
  \begin{subfigure}[t]{0.45\textwidth}
    \centering
    \includegraphics[width=8cm]{../../../../Image/BookV/ting7c.png}
  \end{subfigure}
  \caption{Chiếc loa liên tục xuất hiện ở những nơi khác nhau.}
\end{figure}

``\vie{Thứ này bị ám à!?}'' tôi thở dốc, mắt đảo quanh.

``\vie{Ông ta đang muốn chơi khăm chúng ta đây mà!}'' cô chị bò xuống sàn, nhảy vút lên chiếc loa vừa biến mất ở gần tiêu đề.

``\vie{Cứ như trò đập chuột ấy!}'' cô em gái reo lên, vừa hào hứng vừa thở hổn hển.

Và rồi\dots

*BỘP!*

Suzuki chạm được vào cái loa.

Cái biểu tượng X bên cạnh chiếc loa lại hiện ra, và một tiếng rên dài phát ra từ trong loa: ``\textbf{MmmMMMMMMmmfgh-!!}''

Kaigou-chan ôm bụng cười ngặt nghẽo: ``\vie{Trò này vui đấy!}''

Tôi cũng cười, lần này là cười to thật sự. Có cái gì đó trong trò rượt bắt ngớ ngẩn này khiến bầu không khí bớt nặng nề hơn hẳn, dù chỉ một chút.

Cái loa hiển nhiên sau đó tự bật lại, trở lại chỗ cũ - vị trí ban đầu nơi nó từng rơi xuống. Giọng của Game nghiến lại, đầy tức giận: ``\eng{\textbf{Dừng lại đi!}}''

*RỘP!*

Chưa kịp phản ứng gì, tôi đã thấy nó bị nhốt gọn trong một chiếc hộp gỗ nhỏ có tay cầm, như thể vừa bị tống giam giữa tòa.

\img{ting8}

``\textit{Ông ta nhanh thế}'', tôi nghĩ, vẫn còn chưa hết ngạc nhiên.

``\eng{Tôi nghĩ thế là quá đủ rồi đấy.}'' giọng Game lạnh tanh thốt ra từ hư không.

Kaigou-chan nhăn mặt rõ rệt: ``\eng{Ể\~{}? Bây giờ mới bắt đầu vui mà!}''

``\eng{Xin lỗi nha, tôi ghét phải làm điều này,}'' ông ta tiếp tục, giọng đều đều như vừa giữ lại bình tĩnh, ``\eng{nhưng tại cô không cho tôi lựa chọn nào khác đâu.}''

Tôi thở dài, cố làm dịu tình hình: ``\vie{Ông ấy nói đúng đó. Đừng có phá phách nữa--}''

``\eng{Vậy tôi sẽ làm cách này!}'' cô nàng cắt ngang, nhe răng cười hóm hỉnh rồi lùi lại ba bước.

``\eng{Khoan-!}''  tôi chưa kịp ngăn thì\dots

*RẦM!*

Cô nàng chạy tới và tung cú đá thẳng vào cái hộp gỗ, và chỉ trong một khắc, nó vỡ vụn như trò chơi domino. Từng mảnh gỗ bay tung tóe.

``\eng{Ha ha! Thế nào!}''

Tôi phì cười, còn Suzuki tròn mắt sững sờ như thể vừa xem phim hành động. ``\vie{Onee-sama mạnh ghê!}''

Game gầm gừ đầy tức tối, quay sang về phía tôi, gằn giọng đầy ẩn ý: ``\eng{Em gái song sinh của cậu đúng là bướng thật đấy, người dùng ạ.}''

Tôi sững lại. ``\eng{\dots Chúng tôi chỉ mới gặp thôi,}'' tôi đáp, nhíu mày, ``\eng{và\dots\ ừ, cô ấy bướng bỉnh thật. Tính cô ấy như thế từ khi trước khi gặp rồi.}''

Người bị chúng tôi nhắc tên không hề phủ nhận. Cô nàng khoanh tay, ngẩng mặt thách thức: ``\eng{Bướng thì bướng chứ, ông có cách nào cho chúng tôi rời khỏi đây không?}''

``\eng{Tôi bảo rồi, ở đây chẳng có gì đâu.}''

``\eng{Vậy thì tôi sẽ\dots}''

Cô nàng đưa tay về phía cái loa, định nhấn vào cái nút tắt tiếng lần nữa, nhưng chưa kịp làm gì thì Game hét lớn: ``\eng{\textbf{Đừng nghĩ đến việc đó!}}''

Ngay tức khắc, một khối kim loại dày cộm xuất hiện quanh cái loa như một lớp giáp chắn. Một thứ gì đó nặng nề, xám xịt, và phản chiếu ánh sáng như sắt nguyên chất.

Kaigou-chan khựng lại nửa giây, rồi cất giọng mỉa mai: ``\eng{Ối chà, lại còn cứng đầu à? Vậy để tôi phá tiếp!}''

Tôi lùi một bước: ``\vie{Ấy, đừng! Nó không phải là khúc gỗ đâu!}''

Cô nàng phớt lờ tôi hoàn toàn, lùi về sau như một cầu thủ chuẩn bị sút bóng, lấy đà ra xa, rồi\dots

*RẦMMMMMM!*

Tiếng va chạm vang lên chát chúa. Cô ấy đứng khựng lại, mắt trợn lên, rồi\dots ``\textbf{AAAAAAÁ!!!}''

Cô nàng nhảy lò cò, tay ôm lấy bàn chân vừa đá, mặt tái đi vì đau. Cô ấy dùng đúng là cái chân cô đụng tường lúc mới tới đây, và cũng chính cái chân mà cô dùng để đá vào chữ O ban nãy.

Suzuki đứng gần đó đổ mồ hôi hột: ``\vie{Onee-sama, chị đá đúng chân lúc nãy luôn đó\dots\ Chị ổn chứ?}''

``\vie{Ái da\dots\ đau quá\dots}'' Kaigou-chan rơm rớm nước mắt, ngồi bệt xuống, cố xoa xoa cổ chân.

Tôi chép miệng, ngồi xuống cạnh cô nàng: ``\vie{Tôi đã bảo rồi mà. Ai bảo cứ hốc hách làm gì.}'' Tôi biết rất rõ mọi thứ xảy ra đang đúng theo kịch bản của tựa non-game mà chúng tôi chơi, nhưng tôi không ngờ Kaigou-chan lại hăng máu tới mức như thế.

``Pfft\dots'' giọng Game vang lên, lần này nghe như cố nín cười, nhưng chẳng thành công lắm. ``\eng{Hai người đi mà gọi Siêu nhân đến phá cái đó nha!}''

``\eng{Ể\dots?}'' tôi ngẩng đầu.

``\eng{Để đề phòng thì\dots}'' ông ta tiếp lời, đính thêm một viên ngọc màu xanh lá phát sáng trước một mặt khối kim loại, rồi nói với vẻ thích thú: ``\eng{Một ít kryptonite. Chắc các cô cậu cũng biết ý của tôi là gì, đúng chứ?}''

\img{ting9}

Tôi liếc sang Suzuki, em ấy còn đang mím môi cố nhịn cười, rồi đáp: ``\eng{Nhưng mà Game-san, Siêu nhân có thật đâu?}''

``\eng{Tôi biết, nhưng vì chúng ta đang ở trong một thế giới game-}''

``\eng{Không-phải-trò-chơi,}'' cô chị ngắt lời, xoa xoa cổ chân, mặt nhăn nhó. ``Owie\dots''

``\eng{À ừ. Tôi quên mất.}'' Game đằng hắng, rồi tiếp: ``\eng{Cái gì cũng có thể xảy ra được, phải không?}''

Tôi ngồi xuống, thở ra một hơi thật dài: ``\eng{Cũng có thể lắm chứ\dots}''

\ct

Tôi và Suzuki ngồi xuống cạnh Kaigou-chan. Cô nàng vẫn đang ôm chân, miệng rên khe khẽ, nhưng không khóc, cũng chẳng than thở nhiều.

``\vie{Chị ráng chịu chút nha,}'' cô bé nhẹ nhàng cầm lấy cổ chân chị mình, bắt đầu xoa xoa.

Tôi gật đầu, cũng làm theo, dùng tay phải nắn nhẹ quanh phần cổ chân bị tổn thương, cố không để chạm vào đúng chỗ vừa bị đá trúng.

``\vie{Cũng không đau lắm đâu,} cô nàng thở ra một hơi, mặt không nhăn nhó như hồi trước, ``\vie{chắc chỉ bầm chút thôi. Tôi cũng bị đầy lần như thế nên quen rồi.}''

Tôi ngạc nhiên nhìn cô ấy. Với lực đá mạnh như vậy mà Kaigou-chan chỉ phản ứng chừng đó\dots\ Quả đúng là phi thường.

Chỉ vài phút sau, cô đã ngồi thẳng lại, chống hai tay ra sau và đưa mắt nhìn về phía khối kim loại xám xịt vẫn sừng sững đứng trước mặt.

``\eng{Vừa nãy Game bảo ta phải gọi ai để đi phá đấy nhỉ?}'' cô chị lớn cười, lém lỉnh hất cằm về phía cái khối kia.

Tôi nhún vai đáp: ``\vie{Si*u Nh*n. Tôi nghĩ ông ta nói thế chả khác nào bảo ta bắc thang lên hỏi ông trời cả. Cái khối kia trông nặng bỏ bố ra kìa.}''

Kaigou-chan cười khẩy, lắc đầu nhẹ: ``\vie{Thế thì tại sao ta phải gọi gã ta trong khi đã có sẵn một Si*u Nh*n ở đây rồi? Lại còn miễn nhiễm với cái viên xanh kia nữa?}''

Cô vừa nói xong, tôi mới chợt hiểu ý. ``\vie{Khoan, cô định-}''

Chưa kịp nói hết câu, cô nàng đã bước tới trước, hai chân đứng vững, cúi người ôm lấy cái khối kim loại đồ sộ ấy và rồi, nhấc bổng nó lên như thể chỉ là một thùng các-tông rỗng.

``\vie{C-Cái quái gì\dots!?}'' tôi thốt lên, trố mắt kinh ngạc.

Cô nàng quay đầu, hai tay nhấp nhổm khối kim loại ấy như đang đong đếm: ``\vie{Chắc cũng chỉ cỡ\dots\ năm trăm cân chứ mấy. Hồi trước tôi nâng mấy thứ nặng hơn rồi.}''

``\vie{N-Năm trăm cân!?}''

Cô em gái của cô ấy cười mỉm, nói với tông giọng như thể đây là chuyện thường ngày ở xã\footnote{Câu thành ngữ này vốn dĩ được cải biên lại từ ``chuyện thường ngày ở huyện,'' với lý do là Việt Nam đã loại bỏ cấp huyện trong cuộc cải cách Sáp nhập tỉnh thành từ 01/07/2025.}: ``\vie{Onee-sama còn từng nâng đồ lên đến hàng tấn kia mà.}''

``\vie{\textbf{C-Cả tấn ư!?}}'' tôi càng ngỡ ngàng hơn.

Kaigou phẩy tay, như thể số đó chẳng là gì cả: ``\vie{Thế vẫn chưa là gì đâu. Còn có lần tôi nâng cả một chiếc xe tải hàng chục tấn, chỉ vì nó đậu chắn đường tôi chở hàng hồi ở thế giới cũ thôi.}''

``\vie{\textbf{Cái gì cơ!?}}'' môi tôi há hốc, người tôi đứng cứng ngắc vì sốc trước thông tin vừa nghe. Cả chục tấn, không phải năm trăm hay một tấn đơn thuần, mà là mười vài tấn của một chiếc xe tải đầy hàng? Sức mạnh đó đã vượt quá giới hạn của bất kỳ con người bình thường nào, thậm chí vượt cả trí tưởng tượng của tôi. Không lẽ nào cô ấy là ``Si*u Nh*n'' phiên bản đời thực!?

Nhưng rồi, tôi bỗng nhớ lại lúc hai chúng tôi chạm mặt nhau lần đầu. Cú đấm trời giáng tới mức nứt sàn ở The Void. Cú đấm khiến tôi suýt nữa đã bỏ mạng đầy oan uổng đó.

Giờ thì mọi chuyện đã hợp lý rồi. Nhưng\dots\ tại sao?

``\vie{Có thể anh sẽ không tin, Onii-sama ạ, nhưng mà để em kể cho anh câu chuyện này\dots}'' Suzuki nhìn cô chị nâng chiếc hộp kim loại không mảy may gì, tay chống cằm. ``\vie{Hồi Onee-sama còn nhỏ, lúc bố mẹ bọn em còn sống, chị bị một trận ốm nặng lắm. Thập tử nhất sinh luôn, tới mức bác sĩ còn bảo là không thể chữa trị được.}''

Cô bé hít một hơi sâu, rồi tiếp tục: ``\vie{Nhưng rồi, may mắn đã mỉm cười với gia đình em. Onee-sama từ một người nằm liệt giường, chị ấy bỗng chốc khỏe mạnh, đến độ các bác sĩ còn không tin nổi vào mắt mình.}''

Tôi gật đầu. ``\vie{Vậy chuyện đó thì liên quan gì đến việc chị của em khỏe đến mức đó?}''

Cô bé nhìn về phía cô chị, nở một nụ cười hơi trầm tĩnh: ``\vie{\dots Sau trận đó, bác sĩ bảo là não chị bị ảnh hưởng nhẹ. Nhưng không phải kiểu giảm trí nhớ hay gì đâu, mà là bộ giới hạn não bộ của chị ấy bị hỏng ạ.}''

``\vie{V-Vậy có nghĩa là\dots}''

Suzuki gật đầu. ``\vie{Con người chúng ta thường chỉ dùng một phần nhỏ sức mạnh để đảm bảo an toàn, thì chị ấy lại có thể dùng gần hết, lúc nào cũng đưa chị ấy vào rủi ro.}''

Tôi nghe mà như nuốt trúng cục đá. ``\vie{Có nghĩa là, sức người thường thì chỉ nâng được vài chục cân, nhưng cô ấy có thể nâng cả tấn mà không biết là mình đang phá định luật vật lý và sinh học thông thường?}''

``\vie{Vâng. Ban đầu chị ấy cũng bất ngờ lắm, nhưng rồi quen dần. Bố mẹ em thì vừa lo, vừa tự hào đấy ạ.}'' Cô em gái cúi đầu, ánh mắt lặng đi một thoáng như vừa nhắc tới một điều gì đó quá khứ mà cả hai chị em không thể nào tin được.

Kaigou-chan không nói gì. Cô quay mặt sang chỗ khác, nhưng tôi nhận ra hai bàn tay cô nắm lại rất nhẹ, dường như đang ngượng ngùng với câu chuyện đó.

Tôi khẽ mỉm cười: ``\textit{Vậy là mọi chuyện bắt đầu từ đó à.}''

``\vie{\textbf{Được rồi! Xem tôi đây!}}'' cô nàng lấy đà, thốt lên đầy tự tin.

Thế rồi, cô ném cái khối kim loại nặng trịch lên không trung. Nó xoay vòng mấy vòng rồi *RẦM!* nó đập mạnh xuống đất, rung chuyển cả không gian. Tôi phải giơ tay lên che mặt vì bụi đất tung mù mịt.

Dòng chữ tiêu đề phía trên bị chấn động, các ô chữ bị nhấc bổng lên rồi rơi rụng, có vài ô nứt ra và vỡ toang, để lộ một thứ gì đó phía sau.

\img{ting10}

``\eng{Đến thế này thì chịu rồi,}'' Game bật tiếng thở dài nặng nề. ``\eng{Cứ phá phách đi, tôi cũng chẳng quan tâm nữa.}''

Tôi, Kaigou và Suzuki đồng loạt nhìn lên chỗ dòng tiêu đề vừa bị rơi xuống. Có gì đó ẩn sau lớp chữ cái ấy. Một thứ gì đó mà ngay cả ông ta cũng đang giả vờ như không thấy.

\pov{Hikawa Kaigou}

Hai con vít đang bắt chặt, gắn chặt một cái nắp đen kịt trông như có gì đó bị giấu sau lớp tường.

Suzu bước đến gần, nghiêng đầu: ``\vie{Nhìn hai cái đó\dots\ giống như kiểu che cái gì lại ấy.}''

Kuon-kun thì đứng khoanh tay, mắt nheo lại. ``\vie{Tôi nghĩ tôi từng thấy kiểu vít này ở chỗ thùng máy tính của Botan-chan. Cái đầu vít hình sao năm cánh. Từ từ\dots}''

Cậu ta đặt ba-lô xuống đất, mở toang ra lục lục một hồi.

``\textbf{À ha!}'' cậu giơ cao một cây tua-vít đầu sao. ``\vie{Tôi quên chưa bỏ lại cây này. Lúc sửa máy cho em ấy, định mang về để tra keo mới mà lại quên.}

Tôi rướn mày, mặt nhoẻn đầy thích thú và nhẹ nhõm, rồi cúi người xuống. ``\vie{Thế thì lên đây.}''

Cậu ta chớp mắt: ``\vie{Gì cơ?}''

``\vie{Tôi cõng cậu lên. Cái này cao lắm, tôi với không tới. Nhưng tôi thì khỏe, còn cậu thì nhẹ. Win-win rồi còn gì.}''

``\vie{Ờm, nhưng mà\dots}''

``\vie{Yên tâm, tôi không vật cậu xuống đâu,}'' tôi lườm. ``\vie{Lên đi.}''

Không còn cách nào khác, Kuon-kun đành trèo lên lưng tôi. Tôi đứng thẳng dậy, cậu ta hơi chới với một chút nhưng giữ thăng bằng khá tốt. Suzu thì đứng kế bên, tay cầm chữ cái nhặt được từ tiêu đề rớt ra, vừa nhìn vừa hồi hộp.

``\eng{\dots Hai người nghĩ hai người đang định làm gì thế, những Người dùng?}'' giọng của Game lại vang lên, lần này không giận dữ mà như kiểu lo lắng.

Cậu ta không trả lời, chỉ cắm tua-vít vào con vít thứ nhất và bắt đầu vặn.

*KÉT* *KÉT*

Từng vòng tua quay ra, vít long ra khỏi tường, rơi ``tách'' xuống sàn. Không cần đợi cậu ta ra lệnh, tôi di chuyển sang bên phải, hai tay nắm lấy chân cậu thật chắc, đi tới chiếc thứ hai.

*KÉT* *KÉT* *CẠCH!*

``\textbf{Ah!}''

Chiếc nắp đen bất ngờ rơi phịch xuống, suýt nữa thì trúng đầu Suzuki. Cô bé hốt hoảng lùi lại mấy bước, mặt tái mét, ôm lấy ngực: ``\vie{H-Hú hồn con chim én\dots}''

Tôi vội cúi thấp người để Kuon-kun nhảy xuống, trong khi cậu ta cũng đang hoảng vì cú rơi bất ngờ ấy.

Phía sau nắp tường\dots\ hiện ra một bảng điện tử với bốn ô chữ đang chờ được điền, nằm ngay dưới dòng ``THERE IS A''. Các chữ cái rời rạc rơi vãi khắp nơi, tôi thấy vài cái nằm dưới sàn ngoài chữ O vừa bị rơi rụng ra.

\img{ting11}

Ngay khi chúng tôi chưa kịp định hình, giọng Game thét lên: ``\textbf{\eng{Không! Đừng nhìn! Thứ này là riêng tư đấy!!}}''

Tôi liếc qua Suzu, con bé đang nhón chân tới gần cái bảng, mắt sáng rực như đứa trẻ thấy trò chơi mới.

Kuon-kun thì khoanh tay, gật gù: ``\vie{Vậy là đằng sau cái tiêu đề chính là chỗ quan trọng nhất rồi.}''

``\eng{Mấy người rồi sẽ hối hận vì làm điều này đấy!}'' ông ta vang lên, lần này không còn là lo lắng, mà là giận dữ và có phần cay cú.

Tôi chống nạnh, nhướng mày khiêu khích: ``\eng{Miễn là tôi tìm được cách rời khỏi đây, tôi sẽ không hối hận đâu!}'' tôi bồi thêm bằng nụ cười mỉa mai. ``\eng{Cùng lắm thì ông bị khui vài bí mật nho nhỏ thôi mà\~{}}''

``\eng{Cô\dots\ cô đúng là loại con gái bốc đồng, thô lỗ, đáng ghét, bạo lực, khó ưa\dots}''

``\eng{Tiếp đi, tiếp đi,}'' tôi khoái chí ngắt lời. ``\eng{Xem ông còn từ gì nữa không. Mấy cái lời đó tôi nghe chán chê rồi.}''

``\eng{Và\dots\ và\dots\ đúng là một con đàn bà lẳng lơ!}''

\dots Hửm?

Tôi khựng lại trong chôc lát, rồi lập tức bẻ ngược cổ tay, ngửa mặt lên giả vờ rưng rức khóc: ``\vie{Suzu\~{} Kuon-kun\~{} Ông ta bảo tôi là gái bạo lực, khó ưa với bướng bỉnh kìa\~{}}''

Cô em gái của tôi đứng kế bên, tay che miệng cười khúc khích, như thể vừa nghe hai chúng tôi cãi tay đôi nhau: ``\vie{Chị thì lúc nào chả thế. Nhiều khi chị còn làm quá lên đấy chứ.}''

Cậu ta thì thở dài, gật gù không chút do dự: ``\vie{Ông ta nói đúng quá rồi còn kêu cái gì nữa. Cô còn suýt giết tôi ở The Void lúc trước thế.}''

Tôi quay lại, lườm Kuon-kun bằng ánh mắt sắc nhọn: ``\vie{Phản bội nhanh thế à?}''

``\vie{Có gì thì thành thật như thế thôi. Thẳng thắn thành thật thường thua thiệt, nhưng sự thật là sự thật.}''

``\vie{Thành thật kiểu đó thì\dots\ ăn dép bây giờ.}'' Tôi dậm chân, phồng má giận dỗi, nhưng chỉ dậm nhẹm chứ đau chân lần nữa thì chẳng oai chút nào. Tôi quay lại ném ánh mắt giận dỗi về phía Game, nhưng ông ta giờ im bặt như đang nhai cục tức trong cổ họng.

Trong lúc đó, Kuon-kun lại bắt đầu lẩm nhẩm, mắt liếc qua đống chữ cái rơi vãi dưới sàn. Một chữ T, một chữ R, hai chữ E, một chữ G, một chữ A, một chữ M và một chữ O.

Suzu cũng cúi xuống nhặt lấy mấy chữ, lật qua lật lại như đang chơi Scrabble.

Tôi lại gần, nheo mắt: ``\vie{Có vẻ đủ để tạo thành một vài từ.}''

Kuon gật đầu: ``\vie{Ừ. Mấy từ tiếng Anh có bốn chữ cái đơn giản. Có mấy từ mà tôi nghĩ được này.}'' Cậu liếc mắt về cô em gái của tôi, rồi cùng nhau gật đầu như có cùng tư tưởng.

\pov{Hikawa Kuon}

Suzuki vỗ hai tay một cái, rồi reo lên như thể vừa chợt nhớ điều gì: ``\vie{A! Game-san vừa bảo là ở đây không có trò chơi nào cả, đúng không? Vậy\dots\ nếu mình biến nơi này thành nơi CÓ trò chơi thì sao?}''

Tôi quay sang nhìn cô bé. Cô ấy nói đúng. Quá đúng. ``\vie{Hợp lý đấy, Suzuki,}'' tôi cười nhẹ.

Kaigou-chan cũng không chậm hơn nửa giây, giơ nắm đấm lên cao: ``\vie{Đập tan tuyên bố của ông ta luôn!}''

``\vie{Ừ. Làm nghịch lý tồn tại.}''

Tôi cúi xuống, chắp hai tay ngang lưng. Suzuki hiểu ngay ý, nhẹ nhàng trèo lên lưng tôi với sự phấn khích không giấu được trên khuôn mặt. Trong lúc đó, Kaigou-chan nhanh chóng nhặt lấy từng ô chữ rồi chuyền lên.

``G\dots''

Cô bé trên lưng đón lấy, gắp gọn vào ô đầu tiên.

``A\dots''

Tôi giữ vững tư thế. Cô bé đặt tiếp vào vị trí thứ hai, rồi quay lại giơ tay.

``M.''

Cô chị lớn cười tươi, đặt ô chữ tiếp theo vào lòng bàn tay cô em gái mình, người chậm rãi lắp nó vào đúng vị trí, rồi cúi đầu một chút như thể để tập trung thật kỹ.

``\vie{Chữ cuối nhé, E!}''

Ngay khoảnh khắc đó, một luồng điện chạy dọc sống lưng tôi. Tôi cảm tưởng có điều gì đó trong không khí vừa thay đổi, hoặc chính xác hơn, ngay từ khoảnh khắc mà chúng tôi nghĩ ra trò này.

Giọng Game bất ngờ thốt lên, không còn điềm tĩnh hay mỉa mai như trước nữa, mà là bấn loạn: ``\eng{M-Mấy người\dots\ mấy người đứng có nghĩ là\dots!}''

Nhưng chúng tôi không thèm để tâm. Suzuki dứt khoát ấn ô chữ cuối cùng vào.

*TẠCH!*

Tiếng động phát ra khi bốn chữ cái ăn khớp hoàn hảo với khung chữ.

Ngay lập tức, mặt đất rung chuyển. Một cơn động đất khác lại bắt đầu, mạnh hơn lần trước, như thể toàn bộ cái thế giới kỳ lạ này đang gào thét phản đối.

``\eng{\textbf{Không, không, không, không!}}'' ông ta rú lên, hoàn toàn hoảng loạn.

Một mũi tên nhấp nháy màu xanh trắng hiện ra bên mé trái màn hình. Đồng thời, một vài mảng hồng rách toạc khỏi bức tường tối om, như thể cái vỏ bọc dối trá nào đó vừa bị xé rách.

Tôi quay sang Suzuki. Cô bé vẫn còn đang bám trên lưng tôi, miệng nở nụ cười nhỏ.

Kaigou-chan khoanh tay, gật gù: ``\eng{Đấy. Nơi không có trò chơi bây giờ đã có trò chơi rồi. Vấn đề chỉ là ông có đủ can đảm để đối mặt không, Game-san!}''

Chúng tôi chạy theo mũi tên xanh lá vừa xuất hiện, băng qua một bức tường vừa hé mở. Cứ như thể có một cánh cửa vô hình đã được mở ra sau khi từ được lắp xong.

\ct

Bên trong là một căn phòng mới, sáng hơn, với một bảng điều khiển lớn đặt giữa trung tâm. Ngay khi tôi định bước tới, thì ánh sáng nhấp nháy lên một lần nữa, và một màn hình trò chơi hiện ra trước mặt.

Tôi chớp mắt. Suzuki reo lên: ``\vie{Onii-sama, cái này nhìn quen lắm!}''

Kaigou-chan trố mắt ngạc nhiên: ``\vie{Đây là\dots\ Brick Breaker sao?}''

Một trò cổ điển với một bàn chơi màu hồng, với các hàng gạch màu tím sáng lấp lánh trải dài phía trên, và một quả bóng nằm trên thanh phóng phía dưới. Mọi thứ đều mang sắc hồng neon, như thể đây là một chiều không gian khác hoàn toàn với nơi tôi từng biết.

\img{ting12}

Tôi bật cười thích thú: ``\vie{Trời ạ, tôi nhớ trò này. Hồi đó còn chơi trên mấy cái điện thoại cục gạch nữa cơ.}''

Cô chị gật đầu: ``\vie{Đơn giản mà vui. Cho tôi chơi với nhé?}''

Tôi chỉ vào màn hình: ``\vie{Mỗi người một bên. Nửa màn hình trái cho tôi, phải cho cô. Suzuki điều bóng ở giữa.}''

``Vâng ạ\~{}'' Suzuki hào hứng.

Ngay khi tôi chạm vào thanh điều khiển, quả bóng bật khỏi thanh phóng, bay lên, đập vào một viên gạch ở hàng dưới cùng rồi nảy xuống. Viên gạch biến mất với tiếng *BỤP* nhẹ.

Ngay lập tức, Game rú lên: ``\eng{\textbf{Không! Không được! Rời khỏi đó ngay! Tôi cấm cô cậu đụng đến nó!}}''

Kaigou-chan giả vờ bĩu môi, phụng phụ: ``\eng{Ể\~{}? Nhưng mà chán òm à\~{} Ông có trò chơi mà không cho người ta chơi sao?}''

``\eng{Ông giấu kỹ quá, chúng tôi chẳng biết gì cả,}'' tôi cũng theo đà thế mỉa mai.

Ông ta hoảng loạn, lắp bắp chống chế: ``\eng{T-Tôi cấm! Nếu các người chơi\dots\ chúng sẽ tới!}''

Tôi khựng tay một chút. Tôi biết ông ta đang nói tới cái gì. Nhưng chúng tôi vẫn đồng thanh hỏi: ``\eng{Chúng? Là ai?}''

Game đáp, run rẩy: ``\eng{Mấy người thực sự không muốn biết đâu. Nếu chúng xuất hiện, tôi tiêu chắc\dots}''

Tôi liếc nhìn Suzuki. Dù hơi lo lắng, nhưng cô bé vẫn bình tĩnh điều khiển bóng. Tôi gật đầu trấn an.

``\vie{Kệ đi. Trò chơi là trò chơi. Giải trí chút cũng đâu có hại gì.}''

Chúng tôi tiếp tục, mặc kệ những tiếng rầy la từ giọng nói đặc sệt xứ Pháp kia. Quả bóng đập thêm vài viên gạch nữa. Kaigou-chan thì hí hửng vỗ tay như đứa trẻ, còn Suzuki thì điều chỉnh tốc độ rất chuẩn.

Ông ta hét lớn, gần như mất kiên nhẫn: ``\eng{\textbf{Tôi cảnh cáo rồi đấy! Đừng để tôi phải\dots}}''

Tôi vừa định nói gì thì quả bóng bật lên tận hàng gạch trên cùng, rồi nảy ngược lại, phá liên tiếp mấy ô liền nhau. Cô chị lớn reo lên: ``\vie{Hay quá! Combo luôn!}''

Game càng tức tối hơn, ông ta rít lên: ``\eng{\textbf{Đừng phá gạch nữa!}}''

Bất chợt, các viên gạch đột nhiên tự dịch chuyển. Chúng tránh quả bóng như tránh tà, trượt sang bên trái hoặc phải theo kiểu rất phi tự nhiên.

\img{ting13}

``\eng{\textbf{Ê! Chơi ăn gian!}}'' cả ba chúng tôi hét lên, nhưng vẫn không giấu nổi niềm thích thú khi chúng tôi vẫn phá được thêm vài viên nữa.

``\vie{Điều bóng sang phải một chút!}'' tôi hô lên, điều khiển bệ lệnh sang trái.

``\vie{OK liền\~{}}'' Kaigou-chan đáp lại. Dù rằng ông ta cố tình làm trò khó hơn, nó vẫn khiến chúng tôi kích thích hơn.

Tới đây, giọng nói trầm kia tỏ vẻ giận đến sủi bọt trong lúc ba chúng tôi vẫn cười đùa: ``\eng{Tôi phải tìm cách ngăn các người lại\dots}''

Một tiếng ``Hừm'' thoáng thốt lên, rồi bất chợt ông ta bật ra, như thể vừa nghĩ ra một ý tưởng hay: ``\eng{À, đúng rồi!}''

Ngay khoảnh khắc quả bóng vừa chạm một viên gạch khác, nó nổ tung. Một đám khói pixel màu tím bung ra.

Màn hình im lặng. Không còn bóng.

``\vie{\dots Hết bóng rồi,}'' Kaigou thoáng đơ một hồi. Tôi lắc đầu, cười mếu: ``\vie{Trời ạ, đang hay mà\dots}''

Trông thấy vẻ mặt thất vọng của chúng tôi, Game đầy vẻ đắc thắng: ``\eng{Giờ thì không có bóng nữa, khỏi chơi!}''

Kaigou-chan phồng má lên: ``\eng{Đúng là tên đê tiện\dots}'' còn Suzuki thì cúi đầu buồn buồn, như thể bị tước mất món đồ chơi yêu thích.

Tôi nghe thấy tiếng ông ta cười. Một tiếng cười nhỏ, đầy mưu mô, đắc ý như một đứa trẻ vừa giành lại được món đồ chơi khỏi tay người khác.

Hoặc ít nhất\dots\ ông ta nghĩ vậy.

Đúng lúc này, cô em gái lên tiếng: ``\vie{Này anh chị ơi! Nhìn kìa!}''

Tôi nhìn theo tay cô bé chỉ. Ở góc trái màn hình, ngay dưới dòng điểm, hiện ra một biểu tượng viên cầu màu tím lấp lánh, ngay bên cạnh dòng chữ ``$\times 2$''.

Cô chị lớn nheo mắt: ``\vie{Đó là quả bóng còn lại phải không?}''

Cô nàng khẽ mỉm cười rồi gọi lớn: ``\eng{Game-san\~{} Nếu trò chơi này không còn bóng  thì cái kia là cái gì thế?}'' rồi chỉ tay về biểu tượng số mạng còn lại.

Một khoảng im lặng đáng ngờ.

Thế rồi, Game ấp úng: ``\eng{C-Cái đó là\dots\ hình tượng trưng thôi! Đúng vậy, chỉ là biểu tượng thôi!}''

Tôi nhướng mày: ``\eng{Thế sao nó trông y chang giống quả bóng vậy?}''

Ông ta tiếp tục chối bay chối biến: ``\eng{T-Tựa game nào mà chẳng có biểu tượng như thế!}''

Kaigou-chan khoanh tay, nhìn tôi và Suzuki một lượt rồi gọi lại như hội ý: ``\eng{Thôi được.}\ \vie{Suzu, Kuon-kun. Lên nào.}''

Tôi khẽ gật đầu, cúi xuống để cô em nhỏ leo lên lưng. Cô bé hơi run một chút, nhưng rồi cũng ổn định vị trí. Trong khi đó, người chị lớn dùng hai tay nâng tôi lên cao, khiến cả ba thành một tháp người chồng chéo trong không trung. Làm tôi nhớ đến cách mà Coco-chan chồng Watame-chan và Luna-chan lên để che mưa\footnote{Xem lại \flaregothic{ÉPISODE 112}, quyển số Chín.}.

Tôi chạm tay tới biểu tượng.

Chạm nhẹ.

Một tiếng *TINH* phát ra, quả bóng thứ hai bay khỏi góc màn hình, chạm xuống nền, lăn vài vòng rồi nằm yên tại chỗ xuất phát.

Tôi cẩn thận đặt quả bóng đó lên đúng thanh phóng ở vị trí đủ hiểm để phá gạch theo đường chéo. Suzuki nhỏ giọng: ``\vie{\hl{Onii-sama và Onee-sama có chắc là chỗ này không?}}''

Tôi gật đầu. ``\vie{\hl{Cứ tin ở anh chị. Anh chơi nhiều nên mới biết chứ.}}''

Tôi thả bóng.

*BỘP! BỘP BỘP BỘP BỘP!*

Quả bóng đập vào một góc bảng chơi, phá liên tiếp năm viên gạch theo đường xiên.

``\vie{Thế mà nó phá được kìa!}'' tôi cười lớn.

Kaigou-chan cười khẩy: ``\vie{Tôi tinh mà. Biết chứ?}''

Nhìn chúng tôi tìm được cách phá luật để tiếp tục trò chơi, Game chỉ có thể sững sờ mà không nói nổi một lời nào.

Sau một hồi im lặng, ông ta thở dài bó tay chịu trận: ``\eng{Được rồi. Chả còn vui nữa đâu.}''

Ngay khi lời đó dứt, toàn bộ khung cảnh lại rung lên dữ dội.

Màn hình bắt đầu nhòe đi, như thể bị trầy xước, méo hình. Quả bóng bỗng nặng nề hơn, mỗi lần nảy lên như chậm lại.''

``\vie{Cảm giác này\dots\ kỳ quặc thật.}''

Một đường sọc đỏ-xanh-lục xuất hiện từ góc bên trái màn hình. Rồi thêm một đường nữa. Rồi một mảng. Những khối ô nhiễu RGB bắt đầu lan ra, như bị lỗi shader hoặc xé hình khi chưa bật V-Sync.

\img{ting14}

Suzuki rùng mình, ôm lấy tôi từ phía sau. Tôi cũng toát mồ hôi, dù biết rất rõ chúng là gì ngay từ đầu.

Kaigou-chan khựng tay lại: ``\vie{Cái\dots\ Cái quái gì vậy?}''

Game thốt lên, tông giọng hiện rõ vẻ sợ hãi: ``\eng{Ôi không\dots\ chúng đến rồi!}

Cả ba chúng tôi nhìn xung quanh bảng chơi, tay vẫn đẩy qua đẩy lại bệ phóng: ``\eng{Chúng!?}''

``\eng{Lỗi! Là lỗi đấy! Chúng ở khắp mọi nơi!!}''

``\eng{Thế giờ sao!? Chúng tôi phải dừng lại bằng cách nào?}'' cô chị lớn luống cuống điều khiển bệ phóng.

Giọng nói kia hoảng loạn: ``\eng{Không biết! Hai người phải tìm cách nào đó để diệt chúng đi!}''

``\eng{Nhưng mà\dots}''  tôi đáp, ``\eng{tôi không dừng game được!}''

Game kêu lên, tỏ vẻ bất lực tột độ: ``\eng{Đúng là thảm họa! Làm sao tôi có thể trở thành một tựa game thành công với đống lỗi thế này!?}''

Tôi đổ mồ hôi hột, nhưng vẫn giữ bóng trong tầm kiểm soát: ``\eng{Đó là điều ông quan tâm sao!?}''

Ông ta gào lên, dáng điệu như đang cố động não của mình: ``\eng{Nghĩ đi, nghĩ đi nào!}''

Tôi nghiêng người tránh một cụm nhiễu vừa xuất hiện, bắt đầu đề xuất: ``\eng{Sửa lại code chăng?}''

Kaigou-chan vẫn cứng tay giữ vững thanh phóng: ``\eng{Khởi động lại không-phải-trò-chơi?}''

``\eng{Nghe hay đấy\dots\ nhưng\dots}'' Game tỏ vẻ chần chừ.

Tôi nghiến răng: ``\eng{Nhưng sao?}''

Giọng ông ấy đột ngột trùng xuống: ``\eng{Tôi\dots\ mất quyền điều khiển rồi.}''

Cả ba chúng tôi liếc nhìn nhau, khi này mọi chuyện dần đã nghiêm trọng. Mọi chuyện diễn ra y như ở trong tựa non-game đó, và mặc dù chúng tôi biết chuyện gì xảy ra\dots

\dots việc mà nó xảy ra ngay lúc này lại là một trải nghiệm khác. Không đến nỗi kinh hoàng, nhưng vẫn đủ hỗn loạn để chúng tôi rùng mình.

Tới đây, cả hai chúng tôi cùng thốt lên: ``\eng{\dots Bỏ mịa rồi.}''

Những khối màu nhiễu RGB ngày càng nhiều lên, rì rầm bên tai là tiếng xé vỡ của hệ thống, như thể có ai đó đang cố tách rời từng pixel khỏi nhau. Ngay cả ánh sáng quanh tôi cũng bắt đầu không còn là ánh sáng nữa, nó như một lớp màn rung rinh bị lỗi khung hình.

Tôi không nói ra, nhưng\dots\ Tôi cảm thấy e sợ.

Tôi biết Suzuki cũng cảm thấy vậy. Cô bé im bặt, tay nắm chặt lấy tay tôi. Còn Kaigou-chan, dù vẫn giữ dáng điệu điềm nhiên, tôi thấy cô ấy khẽ nhíu mày, có lẽ cũng lo.

Không ai nói gì thêm. Chúng tôi chỉ biết: phải phá nốt viên gạch cuối cùng.

Tôi đưa bóng về đúng góc phản hồi lần trước, căn chỉnh thật chuẩn, rồi thả tay.

*BỤP!*

Tiếng đập quen thuộc vang lên, viên gạch cuối cùng vỡ tan, một chiếc cúp màu hồng tím hiện ra lơ lửng giữa không trung. Nó chớp sáng nhẹ, kêu một tiếng *TEN TÈN TEN* như đang gọi mời người chiến thắng bước lên bục danh dự.

\img{ting15}

Mỗi tội, nó ở vị trí rất cao, còn cao hơn cả vị trí mà quả bóng thứ hai đặt.

``\vie{Lại phải dựng tháp người nữa hả\dots}'' tôi thở dài, định leo lên Kaigou-chan lần thứ hai.

Nhưng cô nàng lắc đầu, mặt vẻ nghiêm nghị: ``\vie{Không cần. Để tôi.}''

Cô nàng thản nhiên tháo một bên dép đang mang, xoay xoay nó trên tay vài vòng như vận động cổ tay. Tôi vừa định hỏi cô định làm gì, thì, *VÚT!*

Chiếc dép bay thẳng như phi tiêu, găm trúng chiếc cúp.

*KENG!* *BỤP!*

Chiếc cúp rơi xuống đất, lăn mấy vòng rồi dừng lại ngay chân Suzuki.

``\vie{Onee-sama\dots}'' cô bé chớp mắt.

``\vie{\dots làm như đó là chuyện thường ngày vậy,} tôi tiếp lời.

Chúng tôi vừa định cúi xuống nhặt cúp thì giọng của Game lại vang lên, nghe hoang mang: ``\eng{Còn cách gì không nhỉ\dots\ Làm sao xử lý đống lỗi này đây\dots}''

Tôi ngước lên cao, hỏi với sự nghiêm túc: ``\eng{Ông vẫn đang tìm giải pháp à?}''

``\eng{Còn cách nào khác đâu! Mấy người giúp tôi đi!}'' ông ta rít lên, trông bí bách lắm.

Nhìn thấy đối phương đã rơi vào thế yếu, Kaigou khoanh tay lại, đưa ra một yêu cầu: ``\eng{Thế nếu chúng tôi giúp, ông có thể đưa bọn tôi rời khỏi đây chứ?}'

Giọng Pháp kia lập tức vội vã: ``\eng{Giúp! Tôi sẽ giúp! Tôi hứa!}''

Cô nàng nhoẻn miệng, khi buộc kẻ thù cứng đầu nhất ở thế giới này chịu khất phục. Cô nàng giơ ngón tay út ra, nói như ra lệnh: ``\eng{Nói lời phải giữ lấy lời đấy. Ngoắc tay nào.}''

Một khoảng im lặng. Cả bốn người chúng tôi đứng đơ lại, dưới nền của tiếng pixel đang bị nhiễu và những tiếng bụp bụp.

Tôi hắng giọng phá vỡ sự im bặt bất ngờ đó: ``\vie{Ờm\dots\ Kaigou-chan, ông ta không có tay đâu\dots}''

``\dots Ể?'' cô nàng tròn mắt, như vừa không tin những gì cô vừa nghe.

\dots

\dots

``\~{}\~{}\~{}\~{}\~{}\~{}\~{}\~{}'' cô ấy rít lên, đôi má cô nàng đỏ rực. Tôi nhịn không được nữa, phì cười. Game cũng phá lên: ``\eng{Quê là quê là quê cô nàng quê nhiều\~{}}''

``\eng{T-Tôi biết chứ! Tôi chỉ nói đùa thôi!}'' Kaigou-chan hét lên, rõ ràng là ngượng lắm rồi.

Suzuki cười khúc khích sau lưng tôi, nắc nẻ: ``\vie{Onee-sama đúng là ngố tàu thật đấy.}''

\ct

``\eng{Thôi được.}'' ông ta lấy lại giọng nghiêm túc. ``\eng{Vậy làm sao để tôi trở thành một tựa game nổi bật giữa một rổ lỗi như thế này đây?}''

Tôi chống cằm suy nghĩ, rồi đưa ra luận điểm: ``\eng{Nếu sửa code không được, viết lại không được, thì chắc chỉ có thể tận dụng những gì còn sót lại trong hệ thống này.}''

Kaigou-chan gật đầu, đưa ra giải pháp: ``\eng{Chẳng hạn như\dots\ thêm gì đó có sẵn?}''

``\eng{Ừm\dots}'' Game im lặng một lát, rồi liệt kê ra. ``\eng{Nếu vậy, tôi nghĩ được ba thứ.}''

``\eng{Ba thứ?}''

``\eng{Một: những nữ anh hùng mạnh mẽ.}''

Tôi bất giác cười phì. Giọng nói trầm đặc kia hỏi: ``\eng{Cậu cười cái gì vậy?}''

``\eng{Không có gì. Câu đó làm tôi nhớ tới Noel-chan với Roboco-chan thôi. Mỗi tội hai em ấy\dots\ PON cực kỳ.}''

``\eng{PON? Mà, Noel? Roboco? Ai vậy?}''

``\eng{Ngốc nghếch. Hai người ấy cũng là đồng nghiệp của tôi.}''

Ông ta thở dài một tiếng, chuyển sang ý tiếp theo: ``\eng{Thứ hai là\dots\ một đống zombie, và máu me.}''

Cô chị lớn lắc đầu: ``\eng{Nhìn xung quanh cùn chẳng có máu nào cả.}''

Cô em gái thì bám chặt hơn vào tôi, mặt mày xanh xao: ``\eng{T-Tôi không thích máu đâu\dots}''

Tôi tiếp lời: ``\eng{Càng không có zombie ở đây. Với cả mô-típ đó cũng cũ mèm rồi.}''

``\eng{Vậy thì\dots\ thứ ba là\dots}''

``\dots'' cả ba chúng tôi nuốt nước bọt, chờ đợi xem Game sẽ tiết lộ điều gì.

``\eng{\dots một con dê.}''

Cả ba chúng tôi đồng thanh: ``\eng{Con dê?}''

``\eng{Tôi cũng không biết vì sao nữa,}'' ông ta thú nhận. ``\eng{Nhưng tôi nghĩ chúng ta nên thử tìm một trong ba thứ tôi vừa nêu xem.}''

Ngay khi ông ta vừa dứt lời, một mũi tên màu xanh lá hiện lên ở mép trái màn hình. Nó chớp chớp như ra hiệu.

``\vie{Đi thôi,}'' Kaigou-chan gật đầu, cầm lại chiếc dép vẫn còn vương đất.

Tôi cầm chiếc cúp hồng - món quà tưởng thưởng từ trò chơi ban nãy - rồi cùng Suzuki đuổi theo sau cô nàng, bước vào vùng sáng mờ đang chờ sẵn.

\ct

\pov{Hikawa Kaigou}

Chúng tôi quay trở lại căn phòng xếp chữ, chỗ mà mọi chuyện bắt đầu.

Dòng chữ ``THERE IS A GAME'' vẫn còn y nguyên trên bảng, lơ lửng trong không gian u tối, lặng thinh chờ một lần nữa được đánh thức. Những chữ cái còn lại vẫn còn ở đó.

Vẫn là chừng ấy. Vẫn còn cơ hội.

``\eng{Tôi nghĩ,} giọng Game vang lên, như đang gợi ý mà cũng như đang thử vận may cuối cùng của mình, `\eng{lần này\dots\ chúng ta nên thử một trong ba thứ mà tôi đã nói: các nữ anh hùng, zombie, hoặc là\dots\ con dê.}

``\eng{Tôi vẫn không hiểu làm thế nào mà một con dê có thể thu hút người chơi được,}'' Kuon-kun nhíu mày, tỏ vẻ nghi ngờ.

``\eng{Tin tôi đi, tôi cũng không biết tại sao đây này,}'' ông ta đáp lại.

Tôi liếc nhìn đống chữ cái còn lại. Chẳng cần suy nghĩ nhiều. ``GOAT''. Quá dễ. Tôi nghĩ Suzu và Kuon-kun cũng có cùng chung câu trả lời này.

``\vie{Tôi có đáp án rồi,}'' tôi nói, không giấu được chút tự hào.

``\vie{Tôi cũng thế,}'' cậu ta ngẩng lên. Suzu cũng nhìn tôi với sự kiên định.

``\vie{GOAT. Là `con dê'. Chúng ta đang có đủ chữ để ghép.}''

``Onee-sama nói đúng đó,'' con bé gật gù.

Không cần đợi ai chỉ bảo, cả ba chúng tôi lập tức hành động. Tôi đưa từng chữ cái cho Kuon-kun, cậu ấy thì chủ động ngồi xuống để tôi đặt Suzu lên vai như lần trước.

Con bé chồm về phía trước, lần lượt lắp G, rồi O, rồi A, rồi cuối cùng là T đầy nhanh nhẹn.

*XOẠCH* *Ù Ù*

Một lần nữa, cả căn phòng chấn động. Một mũi tên xanh lá lại hiện ra, lần này là ở bên phải màn hình. Nó chớp nháy liên tục, như giục giã chúng tôi bước tiếp.

Tôi là người đi đầu, dẫn hướng cho hai người theo sau.

Ngay sau vạch mũi tên là một đường hầm đá, tối om và ẩm thấp. Không khí ở đây còn vương vấn mùi rêu, mùi đất cũ, mùi ẩm mốc lâu năm trộn với thứ gì đó quen thuộc như từ ký ức game cổ lỗ sĩ.

Chúng tôi chậm rãi tiến vào, cố gắng tránh khỏi những mỏm nhũ đã ở đang cọ sát vào da chúng tôi.

Đường hầm dẫn tới một không gian mở, giống như một hang động trong hầm ngục của những tựa game nhập vai, với vách đá thô ráp và ánh sáng tím mờ lấp loáng từ những khe nứt trong trần.

``\vie{Onee-sama, Onee-sama\dots\ Có cái gì đang treo kìa,}'' Suzu nói nhỏ, tay chỉ về phía trước mắt.

Phía bên kia hang, một cái lồng sắt cao hơn đầu người, bên trong là một con dê màu xanh, trông khá bình thản dù bị nhốt. Bên cạnh cái lồng ấy là một dòng thác nhỏ chảy từ vách đá xuống, rì rào mát lạnh, nước đập vào những hòn sỏi trắng lấp lánh dưới chân thác.

\img{ting16}

``Ơ,'' Game lẩm bẩm. ``\eng{Một con dê kìa! Nhưng mà tại sao nó lại bị khóa trong lồng nhỉ?}''

Kuon-kun gập tay lại: ``\eng{Đúng hơn, tại sao lại có con dê ở đây?}''

``\eng{Tôi cũng chẳng biết,}'' ông ta thừa nhận, nghe như đang nản, lẩm bẩm vế sau, ``\eng{\hl{có lẽ là những gì sót lại khi tôi tạo thế giới này?}}''

Tôi chẳng đợi thêm, bước thẳng lại gần thác nước, ngửa đầu lên và há miệng uống một hơi rõ dài.

Nước ở đây mát lạnh, ngọt lịm và tươihơn hẳn so với nước đóng chai mà chúng tôi uống. Trời ơi, tôi khát tới mức phát điên rồi.

``Onee-sama?'' Suzu tròn mắt ngạc nhiên. Con bé chưa bao giờ uống nước kiểu thế này, mà đã quen với nước đun sôi rồi.

Kuon-kun thì nhăn mặt nhìn tôi: ``\vie{Sao cô lại uống nước đó? Chắc gì nước ấy đã sạch đâu!?}''

Tôi chép miệng, nheo mắt: ``\vie{Tranh luận với ông ta nãy giờ làm tôi khát khô cổ. Giờ mà không uống chắc tôi ngất mất.}''

``Ờm\dots'' cậu ta liếc dòng nước đang chảy ào ào, rồi ngẩng lên hỏi: ``\eng{Nó uống được không vậy?}''

Game đáp lại với giọng nhẹ tênh như gió: ``\eng{Tôi nghĩ là uống được đấy.}''

``\vie{Đó! Ông ta nói là được thì là được thôi. Ông ta đâu có muốn chúng ta chết.}''

Tôi uống thêm vài hớp nữa, tới khi mát lạnh từ đầu tới bụng dưới.

``\vie{Muốn thử không?}'' tôi chìa tay ra với nước được hứng bên trong.

Kuon-kun phẩy tay: ``\vie{Thôi khỏi. Tôi không khát đến vậy.}''

Suzu có vẻ đắn đo một lát, rồi cũng cúi xuống vốc nước lên miệng. ``\vie{Ngon ghê á\dots}''

Tôi gật gù đồng tình: ``\vie{Từ lúc bị lạc vào The Void đến giờ, tôi còn chưa được uống ngụm nào đàng hoàng cả\dots}''

``\eng{Dù sao thì\dots}'' Game đột ngột lên tiếng, kéo chúng tôi về thực tại. ``\eng{Ta cũng tìm được con dê rồi, nhưng dĩ nhiên là nó bị khóa. Chúng ta phải tìm chìa khóa thôi.}''

\il

Câu nói của ông ta trôi tuột qua tai tôi. Bởi vì lúc này, tôi bắt đầu cảm thấy khó chịu.

Mới vài phút trước, nước lạnh mát đến run người, sảng khoái đến tê lưỡi.

Bây giờ, bụng dưới tôi bắt đầu căng tức.

Tôi vừa nhớ ra một điều rất quan trọng, mà đến bây giờ vẫn chưa thể giải quyết được: tôi đang phải nhịn vệ sinh.

Tôi đang mắc tiểu, và giờ thì tôi vừa uống thêm cả nửa dòng thác.

Tệ thật.

Tôi nắm hai tay lại, ấn khẽ vào vùng bụng dưới. Không ổn. Cực kỳ không ổn.

Cả Kuon-kun lẫn Suzu lúc này vẫn còn đang bàn nhau xem có thể tìm chìa khóa ở đâu. Tôi không muốn chen vào.

Tôi cần phải rời khỏi đây. Tôi cần chút thời gian. Một phút. Hai phút.

Tôi lặng lẽ bước lùi lại khỏi cái hang động, quay lưng về phía họ.

Từng bước, tôi giữ chặt hai tay trước bụng, hít thở thật đều, cố không để bản thân quá chú ý tới tiếng nước vẫn rì rào sau lưng.

\pov{-}

Kaigou vừa mới bước ra khỏi hang chưa đầy một phút thì Suzuki nghiêng đầu, hỏi: ``\vie{À-rế? Onee-sama đâu rồi?}''

Kuon quay sang nhìn chỗ cô nàng vừa đứng. Trống trơn. ``\vie{Ủa, nãy còn ở đây mà?}'' cậu lẩm bẩm. ``\vie{Giờ lại chạy đâu mất rồi?}''

Ngay lúc đó, giọng Game vang lên, có vẻ hơi ngập ngừng: ``\eng{Nếu hai người định tìm người bạn đồng hành, cô ấy vừa bước ra ngoài rồi.}''

Cả hai con người cùng quay lại phía cửa hang. Ông ta nói tiếp: ``\eng{Trông cô ta\dots\ có vẻ khá là\dots\ ừm, bối rối. Lúc đi còn hai tay cứ giữ khư khư trước bụng nữa.}''

Cậu con trai khựng lại.

Một thoáng hình ảnh thoáng qua trong đầu cậu: cảnh Kaigou từ chối được cõng lên, bảo là tự làm được, về ánh mắt hơi lảng tránh lúc đó. Và gần đây nhất, cái cách mà cô nàng tu hết cả dòng thác như uống nước lọc mát lạnh sau nhiều ngày đi bộ trong sa mạc, nhưng rồi tỏ vẻ dường như bồn chồn sau khi uống chúng.

Cậu xâu chuỗi lại mọi thứ. Mặt Suzuki vẫn chưa hiểu. Nhưng Kuon thì đã có linh cảm.

``\textit{Mình từng thấy cái biểu hiện đó ở đâu rồi\dots}'' cậu lẩm bẩm, rồi cố nhớ lại một hồi.

Cậu thoáng nhìn về phía cô em gái của người chị, người vẫn đang nghiêng đầu với biểu cảm khó hiểu: ``\textit{À đúng rồi, lúc gặp Suzuki lần đầu.}''

Cậu liếc nhìn cô bé của mình, rồi nhíu mày. ``\textit{\dots Đừng nói là cô ấy đang\dots}''

\begin{center}
  [{\color{red}\textbf{CẢNH BÁO:}} Đoạn văn sau đây có chứa nội dung người lớn và có yếu tố omorashi. Nếu bạn cảm thấy không thoải mái với chủ đề này, hãy cân nhắc trước khi tiếp tục đọc. Chuyển tói trang \pageref{skip} để bỏ qua.]
\end{center}

\pov{Hikawa Kaigou}

Tôi bước ra khỏi hang động, hai tay vẫn giữ khư khư trước bụng, từng bước chân chậm chạp như thể đang đi trên dây thừng. Tiếng nước rì rào từ dòng thác phía sau vẫn vang vọng trong đầu tôi, như một bản nhạc tra tấn không ngừng nghỉ. Bàng quang của tôi đang kêu gào, mỗi cú sốc từ những bước đi làm tôi muốn hét lên.

``\textit{Mình vừa uống cả một dòng thác, Kaigou ơi là Kaigou!}'' tôi tự nguyền rủa bản thân. Ai bảo tôi khát đến mức điên cuồng mà tu ừng ực như thể đó là nước cam ép mát lạnh chứ? Giờ thì tôi trả giá, đúng nghĩa đen rồi đây.

Tôi liếc nhìn Kuon-kun và Suzu, cả hai vẫn đang mải mê bàn bạc về con dê bị nhốt trong lồng sắt. Tôi cố giữ vẻ mặt bình thản, nhưng mỗi lần di chuyển, áp lực dưới bụng lại tăng thêm một bậc. Tôi không thể chen vào cuộc nói chuyện của họ, không phải bây giờ, khi mà tôi đang chiến đấu với chính cơ thể mình.

Trong lúc mò mẫm ra ngoài, mắt đôi đảo quanh tìm kiếm một nơi nào đó để\dots\ giải quyết. Nhưng đây là một thế giới quái gở, không có lấy một bụi cây hay góc khuất để tôi trốn tạm. Chỉ có bóng tối, vách đá, và cái lồng sắt chết tiệt kia.

Đột nhiên, một bóng dáng lướt qua trước mắt tôi, nhanh như một cơn gió. Một con sóc bay, lông xù màu nâu đỏ, đang lượn lờ trên không, đuôi móc một chiếc chìa khóa lấp lánh. Tôi tròn mắt, quên mất cơn đau tức đang hành hạ mình trong một khoảnh khắc.

\img{ting17}

``\textbf{\eng{Nó kìa!!}}'' giọng Game vang lên, khàn khàn và đầy phấn khích, như thể ông ta vừa phát hiện ra kho báu. ``\eng{Con sóc bay đó đang giữ chìa khóa!}''

Hai người theo sau vừa lúc rời khỏi hang đá. Cậu ta nghe ông nói vậy, lập tức lao về phía trước, tay vươn lên cố bắt lấy con sóc, nhưng nó lách mình nhanh như chớp, bay vút lên cao, khuất dạng ở trên cao.

Tôi nghiến răng, cố kìm nén cơn đau để nhảy theo, tay giật lấy một ô chữ cái nằm lăn lóc dưới đất. Tôi nhắm thẳng vào con sóc, hít một hơi sâu, rồi ném mạnh.

*VÚT!*

Ô chữ bay lên, xoay tròn như một chiếc boomerang, nhưng chỉ sượt qua đuôi con sóc trước khi rơi cộp xuống chiếc sàn.

``\vie{Trượt rồi!}'' tôi rít lên, hông rõ là vì tức hay vì cơn đau đang bùng nổ trong bụng. Tôi muốn nhảy lên lần nữa, nhưng mỗi lần rướn người, bàng quang của tôi lại phản đối kịch liệt. Tôi liên tục tự nhủ bản thân không cho ``vòi nước'' bị ``xả ra,'' hai tay nắm chặt vào bụng dưới, cố giữ cho ``con đập'' không vỡ ngay tại đây.

Suzu đã đứng ở bên cạnh tôi từ lúc nào không hay với cặp mắt to tròn đầy bối rối. ``\vie{Onee-sama, chị ổn không? Chị cứ\dots\ đứng kỳ kỳ sao ấy.}'' Giọng cô bé nhỏ xíu, nhưng đủ để khiến tôi giật mình. Chết tiệt, con bé để ý rồi.

Tôi cố nặn ra một nụ cười gượng gạo, lắc đầu: ``\vie{Chị ổn mà, Suzu. Chỉ là\dots\ hơi mệt thôi.}'' Nhưng ánh mắt của con bé vẫn không rời khỏi tôi, như thể đang cố đọc vị điều gì đó.

``\eng{Con sóc bay kia lấy mất chìa khóa rồi,}'' Game lẩm bẩm, giọng đầy ngán ngẩm. ``\eng{Mấy người làm cái gì mà để nó chạy mất thế?}''

Kuon-kun quay lại, thở hổn hển, tay chống hông. ``\eng{Khỉ thật! Giờ chúng ta lấy lại bằng cách nào đây?}'' Cậu ta liếc sang tôi, rồi khựng lại, như thể vừa nhận ra điều gì. ``\vie{Kaigou-chan, cô sao thế? Trông cô như sắp ngất tới nơi rồi ấy.}''

Tôi cắn môi, hai tay vẫn ép chặt vào bụng, mặt nóng bừng. Không thể giấu được nữa rồi. Tôi rón rén giơ một tay lên, như thể đang xin phép giáo viên trong lớp. ``\vie{\hl{T-Tôi\dots}}'' giọng tôi nhỏ xíu, gần như lạc đi vì xấu hổ. ``\vie{\hl{Tôi muốn đi tiểu\dots}}''

Im lặng. Không gian như đông đặc lại. Tôi cảm thấy mặt mình đỏ như cà chua chín, muốn độn thổ ngay lập tức. Tại sao mình lại phải nói ra chứ!? Nhưng cơn đau tức dưới bụng không cho tôi lựa chọn nào khác. Tôi đã cố nhịn từ lúc uống nước ở dòng thác, qua cả màn rượt đuổi cái loa, qua luôn cả việc Suzu ngồi xả lũ đầy vô tư, và giờ thì tôi đang đứng trên bờ vực của sự sụp đổ.

Suzu tròn mắt, che miệng, cố giấu một nụ cười khúc khích. ``\vie{Onee-sama, chị\dots\ muốn đi thật luôn ạ?}'' cô bé nói, giọng vừa ngạc nhiên vừa lo lắng, nhưng vẫn không giấu được chút thích thú trẻ con.

Kuon-kun thì nhíu mày, nhìn tôi từ đầu đến chân, như thể đang ráp nối các mảnh ghép. ``\vie{Tôi đoán được rồi,}'' cậu ta nói, giọng bình tĩnh nhưng không giấu được chút trêu chọc. ``\vie{Cô uống cả dòng thác đó, lại còn nhảy nhót ném dép, đúng không? Bảo sao cô cứ đứng kỳ kỳ từ nãy giờ.}''

Game chen vào, giọng khàn khàn đầy vẻ chế giễu. ``\eng{Cô gái, cô không biết nhịn à? Thế giới này đâu có nhà vệ sinh cho cô đâu. Cố chịu đi!}''

Tôi lườm Game, dù biết ông ta chỉ là một giọng nói vô hình. ``\eng{Ông im đi! Ông có biết cảm giác muốn nổ tung thế này không hả?}'' tôi gần như hét lên, tay vẫn giữ chặt phần hạ bộ, chân nhún nhún như đang đứng trên than hồng. ``\eng{Nãy giờ tôi lỡ uống nhiều nước quá\dots\ với cả đống lỗi pixel lúc nãy làm tôi run muốn chết!}''

Suzu bước đến gần tôi, tay khẽ chạm vào vai tôi. ``\vie{Onee-sama, chị cố lên! Chắc là\dots\ chắc là sẽ có cách mà!}'' giọng cô bé đầy cảm thông, nhưng tôi thấy ánh mắt cô bé lấp lánh chút lo lắng, như thể sợ tôi sẽ ``vỡ đê'' ngay trước mặt.

Cậu con trai thì khoanh tay, lắc đầu mà nói thẳng. ``\vie{Dám cá là nơi này chẳng có nhà vệ sinh đâu. Cô ráng nhịn đi, Kaigou-chan.}''

``\vie{\textbf{Cậu nói dễ lắm!}}'' tôi gào lên, giọng lạc đi vì vừa tức vừa xấu hổ. ``\vie{\textbf{Tôi sắp không nhịn nổi nữa rồi!}}''

Giọng Pháp trầm thở dài, như thể đang ngao ngán. ``\eng{Thôi, kệ cô. Ai bảo cô uống nhiều nước làm gì. Giờ thì tập trung đi, chúng ta cần lấy lại chìa khóa!}''

Tôi muốn phản bác, muốn hét vào cái giọng Pháp đáng ghét đó, nhưng cơn đau tức dưới bụng làm tôi chỉ biết rên lên, hai tay nắm chặt hơn. Kuon và Suzuki quay lại bàn bạc, như thể tôi không tồn tại, tiếp tục nói về con sóc bay và cái chìa khóa. Tôi đứng đó, một mình, vừa giận vừa bất lực, cố kìm nén ``con đập'' đang chực chờ bung ra.

Trong cơn túng quẫn mà tôi đang trải qua, ông ta bỗng lên tiếng, giọng có chút nghiêm túc hơn. ``\eng{Về chuyện con sóc vừa nãy\dots\ Tôi nghĩ cậu nên tìm một cái thang hay gì đó để leo lên.}''

Kuon-kun gật gù, ngầm đồng tình với ông ta. ``\eng{Nhưng ở đây chẳng có thang. Ý ông là gì? Bất cứ thứ gì để trèo lên được?}''

Game cười khẽ, giọng mang chút hoài niệm. ``\eng{Đúng thế. Nó làm tôi nhớ đến hồi thơ ấu\dots\ Tôi từng trèo một cái cây to trong vườn để lấy lại con diều.}''

Tôi trợn mắt, quên mất cơn đau nhất thời: ``\eng{Ông mà cũng có thời thơ ấu á? Ông chỉ là một chương trình thôi mà!}''

``\eng{Cô đừng coi thường,}'' ông ta đáp, giọng hơi hậm hực. ``\eng{Dù là chương trình, tôi cũng có ký ức. Và cái cây đó, nó là một phần của tôi.}''

``\vie{Và cái cây là `TREE',}'' cậu con trai kết luận, như đã biết đáp án từ trước. Cậu ta quay sang tôi, giọng dứt khoát. ``\vie{Kaigou-chan, đưa tôi chữ R, chữ E, và chữ T.}''

Tôi giật mình, cố di chuyển để nhặt các ô chữ cái dưới sàn, nhưng mỗi bước đi là một cực hình. ``\vie{\textbf{Được rồi!}} tôi hét lên, cố che giấu sự run rẩy trong giọng nói. Tôi cúi xuống, tay run run nhặt từng chữ, lẩm bẩm: ``\vie{\hl{Ôi, mót quá\dots}}''

Mỗi lần cúi người, tôi cảm thấy áp lực trong bụng như muốn nổ tung. Tôi đưa từng ô chữ cho câu ta, chậm chạp hơn bình thường, mắt nhòe đi vì đau.

Suzu vẫn đứng bên cạnh, lo lắng nhìn tôi: ``\vie{Onee-sama, chị chậm lại chút đi, không là nguy hiểm đó\dots}'' Cô bé cố giúp, nhặt một chữ E đưa cho tôi, nhưng ánh mắt cô bé đầy lo âu, như thể sợ tôi sẽ ngã quỵ bất cứ lúc nào.

Kuon-kun giờ đã cõng con bé trên lưng, giục tôi: ``\vie{Nhanh lên, Kaigou-chan! Chữ R đâu?}''

Tôi nghiến răng, cố kìm nén cơn đau, đưa chữ R cho cậu ta. Tôi gần như ném ô chữ vào tay cậu ta, chân tôi run rẩy, hai tay lại ấn chặt vào bụng. ``\textit{Chỉ cần xong việc này, mình phải tìm chỗ\dots\ ngay lập tức!}'' tôi tự trấn an bản thân, nhưng trong lòng biết rõ, nơi này chẳng có chỗ nào để tôi ``giải phóng'' cả.

Cậu ta trong lúc đó vừa cõng Suzu lên vai, cẩn thận đặt từng chữ R, E, và T vào bảng chữ, ánh mắt cậu ta tập trung như thể đang giải một bài toán sinh tử. Cô em gái của tôi thì hớn hở, cười tươi khi từng ô chữ khớp vào vị trí, nhưng ánh mắt con bé thỉnh thoảng lại liếc sang tôi, đầy lo lắng. Tôi cố nặn ra một nụ cười gượng gạo, nhưng mặt tôi chắc hẳn đã đỏ bừng như quả táo chín.

Ngay khi chữ cuối cùng được đặt vào, cả mặt đất rung chuyển, mạnh đến mức tôi suýt ngã nhào. Tôi vội chống tay xuống sàn đá, cắn chặt môi để kìm tiếng rên khi áp lực dưới bụng tăng vọt. Một cái cây nhỏ xíu, không hơn gì một cành mầm, đột nhiên trồi lên từ mặt đất ngay dưới chân tôi, lá xanh mướt nhưng yếu ớt như thể chỉ cần một cơn gió là gãy. Tôi tròn mắt, quên mất cơn đau trong một giây, nhưng rồi tiếng rì rào từ dòng thác trong hang động lại kéo tôi về thực tại tàn nhẫn.

\img{ting18a}

``\eng{Ồ, đó là một cái cây nhỏ,}'' Game lên tiếng, giọng khàn khàn pha chút thất vọng. ``\eng{Giá như nó to lớn hơn nhỉ\dots}''

Tôi lườm vào hư không, nơi tôi tưởng tượng giọng nói đó phát ra. ``\eng{Cậu muốn nó lớn hơn? Tưới nước cho nó đi!}'' tôi gắt lên, nhưng ngay lập tức hối hận. Ý nghĩ về nước làm bàng quang tôi co thắt dữ dội, và tôi phải kẹp chặt chân lại, tay ấn mạnh hơn vào bụng.

``\vie{\hl{Ái da\dots{} sắp ra rồi\dots}}'' tôi lẩm bẩm, giọng nhỏ đến mức gần như không ai nghe thấy. Hy vọng rằng mọi thứ cần được làm càng nhanh càng tốt\dots

Suzu quay sang tôi, đôi mắt to tròn đầy lo âu. Tôi muốn trấn an cô bé, nhưng chỉ có thể gật đầu yếu ớt, môi mím chặt. ``\textit{Suzu, chị xin lỗi, nhưng chị đang cố hết sức đây\dots}''

Giọng Pháp tiếp tục tập trung vào câu nói bộc phát của tôi, như thể không để ý đến tình trạng của tôi. ``\eng{Vấn đề là, làm thế nào để đưa nước đến cái cây này? Trước tiên, ta tìm nước ở đâu?}''

Kuon-kun đứng cạnh cái mầm cây, khoanh tay suy nghĩ. ``\eng{Ở trong hang động có dòng thác mà Kaigou-chan uống đến mức\dots\ ừm, muốn おしっこ bây giờ đó,}'' cậu ta liếc tôi, khóe miệng nhếch lên một nụ cười trêu chọc.

Tôi trợn mắt, mặt nóng bừng vì xấu hổ. ``\vie{\textbf{Đừng có khịa tôi, Hikawa Kuon!}}'' tôi gào lên, nhưng giọng lạc đi vì cơn đau đang hành hạ. ``\vie{Cậu thử đứng đây mà nhịn xem, có hét được như tôi không!}'' Tôi muốn lao vào đấm cậu ta một cái, nhưng mỗi cử động lại làm tôi cảm thấy cái ``bể nước ấy'' sắp vỡ.

Suzu che miệng, cố giấu một tiếng cười khúc khích, nhưng ánh mắt của nó vẫn đầy lo lắng. ``\vie{Onee-sama, chị đừng giận Onii-sama mà\dots\ Chị cứ bình tĩnh, em tin chị sẽ ổn thôi!}'' con bé nói, giọng ngọt ngào nhưng run run, như thể đang cố an ủi cả tôi lẫn chính mình.

Game chen vào, giọng đều đều: ``\eng{Rồi, vậy làm sao để hứng nước?}''

Kuon-kun giơ chiếc cúp hồng tím từ trò phá gạch lên, đáp đầy quả quyết: ``\vie{Dùng cái này. Đơn giản mà hiệu quả.}''

``\eng{Hiểu ý đấy,}'' giọng nói trầm kia gật gù, giọng có chút phấn khích. ``\eng{Triển khai đi, rồi lấy lại chiếc chìa khóa!}''

Cậu con trai kia không nói thêm lời nào, lập tức chạy về phía hang động, chiếc cúp lấp lánh trong tay cậu ta. Tôi đứng đó, một mình với cái mầm cây và Suzu, cảm giác bàng quang như đang phình to gấp đôi. Tôi nhìn cái cây, rồi liếc sang con bé, người vẫn đang lo lắng nhìn tôi.

Một ý nghĩ điên rồ lóe lên trong đầu. Tôi khẽ lên tiếng, giọng run run vì xấu hổ: ``\eng{Ờm\dots\ Tôi có thể\dots\ tè vào cái cây này được không?}''

Suzu tròn mắt, há hốc miệng: ``\vie{Onee-sama!? Chị nói gì thế!?}'' Cô bé đỏ mặt, tay vung lên như muốn ngăn tôi lại, như thể không tin tôi dám làm thật.

Game phá lên cười, giọng khàn khàn đầy mỉa mai: ``\eng{Hổ báo lắm mà, cô gái? Giờ lại rúm ró thế sao? Cô nghĩ cái cây này chịu được à?}''

Tôi nghiến răng, mặt nóng ran: ``\eng{T-Tại tôi buồn đi quá chứ bộ!}'' Hai tay vẫn nắm chặt phần hạ bộ, chân nhún nhún liên tục. ``\eng{Rốt cuộc là ông có cho tôi đi không hả?}''

Ông ta buông ra một cái thở dài, giọng bất lực: ``\eng{Tôi e là không được đâu. Cây sẽ chết vì tưới quá nhiều nước đấy.}''

``\eng{Nhưng\dots\ nhưng mà\dots}'' tôi rên lên, giọng gần như van xin, ``\eng{tôi sắp tè ra quần rồi!}''

Cô em gái của tôi bước đến gần hơn, tay khẽ kéo áo tôi: ``\vie{Onee-sama, chị cố thêm chút nữa đi\dots\ Em không muốn thấy chị khổ thế này đâu\dots}'' giọng cô bé nhỏ xíu, đầy lo lắng, và tôi cảm thấy tim mình thắt lại. ``\textit{Suzu, chị xin lỗi, nhưng chị không chịu nổi nữa rồi\dots}''

``\eng{Vậy thì tìm một xó nào đó mà đi đi,}'' ông ta lên tiếng, giọng tỉnh bơ: ``\eng{Tôi bảo rồi, tôi không có mắt. Cứ đi đi, tôi sẽ giả vờ không biết gì.}''

Tôi lườm vào hư không, muốn hét vào mặt cái giọng Pháp đáng ghét đó, nhưng trước khi tôi kịp phản bác, Kuon-kun quay lại, tay cầm chiếc cúp đầy nước. Tiếng nước chảy róc rách từ cúp khi cậu ta đổ xuống cái cây làm tôi muốn phát điên. Tôi kẹp chặt chân, cắn môi đến đau, cố không nhìn vào dòng nước đang thấm vào đất. ``\textit{Không được nghĩ đến nước, Kaigou, không được nghĩ!}''

Nhưng tiếng nước chảy đó như muối xát vào bàng quang tôi, khiến tôi chỉ muốn hét lên.

Cả mặt đất rung chuyển lần nữa, cái cây nhỏ lớn lên nhanh chóng, lá xanh mướt vươn cao hơn. Suzu vỗ tay, reo lên thích thú: ``\vie{Onee-sama, Onii-sama, nó lớn lên kìa!}'' Nhưng ánh mắt con bé vẫn lén nhìn tôi, lo lắng thấy rõ.

\img{ting18b}

``\eng{Xem chừng có vẻ có tác dụng đấy,}'' cậu ta đứng chống hông, liếc nhìn tôi với vẻ quan ngại. ``\vie{Nghiêm túc đấy, Kaigou-chan. Mắc quá thì cứ đi đi. Tôi sẽ không nhìn đâu.}''

``Ư\~{}\~{}\dots'' hai tay tôi nắm chặt hơn, thậm chí tôi còn phải cho ngón tay tôi vào bên trong để bịt trực tiếp ``cổng xả'' lại. Tôi giương ánh mắt hình viên đạn về phía cậu ta, nhưng chỉ đáp lại bằng một cái thở dài chán chường.

Game hét lên, giọng phấn khích: ``\eng{\textbf{Thêm nước nữa đi!}}''

Kuon-kun gật đầu, chạy lại về phía hang động.

Tôi đứng đó, nhìn cái cây đang cao dần, cảm giác bóng đái mình như sắp nổ tung. \textit{Chỉ một chút nữa thôi, Kaigou\dots\ Chỉ một chút nữa\dots} Nhưng tôi biết, tôi không thể cầm cự thêm lâu được nữa. Mặt tôi nóng bừng, mồ hôi lấm tấm trên trán, và tôi biết mình trông thảm hại đến mức nào.

Cậu ta vừa trở lại từ phía hang động không lâu sau đó với chiếc cúp đầy nước lần nữa, liếc sang tôi. ``\vie{Kaigou-chan, cô vẫn đứng được đấy à? Trông cô như sắp ngất lịm đến nơi rồi kìa,}'' cậu ta nói, giọng nửa đùa nửa thật, nhưng ánh mắt cậu ta thoáng chút lo lắng khi thấy tôi run rẩy.

``\vie{Cậu\dots\ im đi!}'' tôi gắt lên, không còn giữ được bình tĩnh được nữa, giọng lạc đi vì xấu hổ và đau đớn. Tôi muốn phản bác thêm, muốn hét vào mặt cậu ta, nhưng một cơn co thắt đột ngột ở bụng dưới làm tôi khựng lại.

Tôi cảm nhận được một giọt nóng hổi, nhỏ xíu, rỉ ra từ lỗ tiểu, thấm qua lớp quần lót cót-tông tôi đang mặc- Không, không được! Tôi hoảng loạn, kẹp chặt chân lại, tay ấn mạnh hơn, nhưng cảm giác ẩm ướt giữa hai đùi làm tim tôi như ngừng đập. Tôi nhìn xuống dưới, thấy một vũng nước vàng đang dần hình thành ngay dưới chân tôi.

Mọi chuyện căng như dây đàn tới mức tới Game cũng không mở lời châm chọc nữa. Hoặc có lẽ, ông ta cũng chẳng buồn cà khịa tôi thêm, khi vấn đề quan trọng hơn còn ở trước mắt. Chỉ là một giọng trầm nhỏ nhẹ như nhắc khéo: ``\eng{Đi thì cứ đi đi. Cơ thể cô không chịu được thêm đâu.}''

Nước mắt bắt đầu cay xè ở khóe mắt. Tôi, Hikawa Kaigou, người lúc nào cũng hổ báo, giờ lại đứng đây, run lẩy bẩy, quần lót ẩm ướt vì không thể kìm nén nổi. Một dòng nước nóng khác rỉ ra, chậm rãi nhưng không thể dừng lại, thấm qua lớp vải quần lót, chảy dọc xuống đùi tôi. Tôi cảm nhận được sự ẩm ướt lan ra, làm chiếc váy của tôi bắt đầu thấm một mảng nhỏ ở giữa, tới mức muốn biến mất khỏi cái thế giới quái gở này.

Và rồi, chuyện gì đến cũng phải đến. Tôi vô thức hét lớn: ``\eng{\textbf{Ra rồi\dots\ Ra rồi!!!}}''

*XÈ\dots*

Từng dòng nước nhỏ giọt kia dần trở thành một dòng lớn trút xuống như thác đổ.

``\vie{Onee-sama!?}'' Suzuki kêu lên, giọng vừa kinh ngạc vừa lo lắng. Cô bé bước đến gần hơn, tay vung lên như muốn làm gì đó, nhưng rồi khựng lại, không biết phải xử lý thế nào.

Tôi không muốn thừa nhận chuyện đang xảy ra lúc này, nhưng cảm giác ẩm ướt giữa hai đùi giờ đã rõ ràng hơn, quần lót của tôi càng lúc càng thấm đẫm hơn, một đường nước khác chảy len lỏi hẳn vào bên trong dép của tôi. Mùi khai nhè nhẹ làm tôi muốn phát điên, và tôi chỉ biết đứng đó, tay che mặt, toàn thân run rẩy vì xấu hổ.

``\eng{Yikes,}'' ông ta bất lực cất lời. ``\eng{Cô ta đi thật à?}''

``\eng{\dots Đúng thế thật,}'' Kuon-kun từ phía sau tôi thở dài. ``\eng{Tôi không muốn tin, nhưng\dots\ sự thật là như thế đấy\dots}''

``\vie{O-Onee-sama à\dots}'' Suzu nhìn tôi với vẻ e ngại, ngỡ ngàng tới độ không thốt nổi một lời nào. Thế rồi, con bé kiên định nhìn về phía cậu con trai: ``\vie{Onii-sama, giúp em với! Đừng để ai thấy Onee-sama như thế này!}''

Kuon gật đầu, dù trông cậu ta vẫn hơi lúng túng. ``\vie{Được rồi, Suzuki. Kaigou-chan, cô cứ\dots\ đứng yên đó. Tôi sẽ đổ nước xong, rồi chúng ta tìm chỗ để\dots\ ừm, xử lý,} cậu quay lại cái cây, cố đổ nốt chỗ nước còn lại, nhưng tôi thấy vai hơi run, như thể vẫn còn đang bàng hoàng với chuyện đang xảy ra.

Cả mặt đất rung chuyển lần nữa, cái cây giờ đã cao vượt đầu chúng tôi, cành lá xòe rộng, che khuất ánh sáng lờ mờ của hang động. Một mũi tên màu xanh lá - trắng hiện ra hướng lên trên, nhấp nháy cùng với hai mũi tên còn lại.

\img{ting18c}

``\eng{Erm\dots\ Đây có phải là lúc để nói\dots\ có thể bắt con sóc kia không nhỉ?}'' giọng Pháp ấy lên tiếng, tỏ vẻ ái ngại cực độ. ``\eng{Tôi ờm\dots\ xin lỗi.}''

Nhưng tôi chẳng còn tâm trí để quan tâm đến cái cây hay chìa khóa nữa. Tôi đứng đó, quần lót ướt sũng, dép quai hậu thì nhớp nháp, nước mắt lăn dài trên má. Suzuki vẫn nắm tay tôi, cố an ủi, còn Kuon-kun thì cố tình nhìn đi chỗ khác, như thể muốn giữ chút thể diện cho tôi.

Tôi thật đúng là thảm hại. Trong lòng tôi bây giờ chỉ còn lại sự bât lực và nhục nhã.

Chiếc váy đỏ của tôi giờ đã ướt nhẹp phần giữa, loang ra một vết sẫm hơn hẳn so với những phần còn lại, lớp quần lót cót-tông bên trong dính chặt vào da, lạnh ngắt và nặng trịch. Mỗi bước nhỏ tôi di chuyển, tôi lại cảm nhận được dòng nước còn sót lại chảy dọc xuống mắt cá chân, thấm vào đôi dép đi dưới chân, để lại mùi khai nhè nhẹ khiến tôi chỉ muốn hét lên vì xấu hổ.

Suzu vẫn nắm chặt tay tôi, đôi mắt to tròn đỏ hoe, như thể cô bé đang cố kìm lệ vì tôi. ``\vie{Onee-sama, không sao đâu\dots\ E-Em sẽ không để ai biết đâu!}'' con bé thốt lên đầy sự cảm thông chính sự an ủi đó lại làm tim tôi thắt lại. ``\textit{Suzu, chị xin lỗi vì đã để em thấy chị thế này\dots}'' Tôi muốn nói, nhưng cổ họng tôi nghẹn ứ, chỉ phát ra một tiếng rên khe khẽ.

Kuon-kun ở gần đó đặt chiếc cúp xuống, quay lại nhìn tôi. Ánh mắt cậu ta lướt qua đôi chân run rẩy của tôi, rồi nhanh chóng nhìn lên, cố tỏ ra bình thường nhưng tôi thấy má cậu ta hơi ửng đỏ.

``\vie{Kaigou-chan, cô\dots\ ừm, cứ bình tĩnh,}'' cậu ta nói, giọng ngập ngừng, như thể đang cố tìm lời an ủi mà không làm tôi thêm xấu hổ. Nhưng tôi biết, cậu ta đã thấy hết: vệt ướt trên váy tôi, những dòng nước chảy xuống từ khe giữa, rồi đến cả dáng vẻ thảm hại của tôi khi đứng run như cầy sấy.

Sự tức giận ngày trước bất chợt bùng lên trong tôi, lấn át cả cảm giác nhục nhã. Tôi nghiến răng, chỉ tay vào khoảng không, nơi tôi tưởng tượng giọng nói đó phát ra.

``\eng{Nếu ông dám nhắc lại chuyện này một lần nữa, tôi thề sẽ tìm cách xóa sạch cái chương trình ngu ngốc kia! Ông không có mắt, nhưng tôi sẽ khiến ông hối hận vì cái miệng đó!}'' tôi bước một bước về phía trước, quên mất tình trạng của mình, và lập tức hối hận khi cảm giác ẩm ướt ở quần lót lại cọ vào da, làm tôi khựng lại, mặt đỏ bừng.

Game thở dài, giọng bớt đi vẻ chế giễu. ``\eng{Được rồi, tôi sẽ không nhắc nữa. Nhưng cô gái à, cô nên làm gì đó với\dots\ ừ, bản thân cô đi.}''

Tôi lườm vào hư không, nhưng trước khi tôi kịp đáp trả, tôi quay sang Kuon-kun, người vẫn đang đứng cạnh cái cây, đứng lùi ra xa từ khi nào.

Ánh mắt cậu ta lảng tránh tôi, như thể cố không nhìn vào thứ đáng xấu hổ ấy. Sự nhục nhã trong tôi lại trỗi dậy, nhưng lần này, nó không phải là cơn giận.

Tôi bước đến gần cậu ấy, chân run run, giọng tôi hạ xuống, gần như van xin. ``\vie{Kuon-kun\dots\ cậu, làm ơn\dots}'' tôi cắn môi, nước mắt lại dâng lên, làm mờ tầm nhìn. ``\eng{Làm ơn quên chuyện này đi. Đừng kể với ai, đừng nhớ gì hết. T-Tôi không muốn ai biết người đồng hành với cậu từng thảm hại thế này\dots!}''

Tôi biết cậu ta là một người khác hẳn so với những gì lũ đàn ông đã thể hiện với tôi, hay những gì mà Game đã làm với tôi ban nãy. Ừ thì, Kuon-kun là một người thẳng thắn, đôi lúc có khi còn thốt ra những lời thô tục tới mức trơ trụi, nhưng cậu ta không phải là người thích lợi dụng những tình huống như thế này để thỏa mãn tâm địa ác độc, mà ngược lại lại nhìn tôi với vẻ lo lắng và e dè.

Cậu ta nhìn tôi, ánh mắt cậu ta dịu lại, rồi mỉm cười nhẹ: ``\vie{Tôi biết rồi, Kaigou-chan. Tôi sẽ không nói gì đâu. Cô cứ yên tâm.}'' Cậu ta đặt tay lên vai tôi, nhẹ nhàng nhưng chắc chắn, như thể muốn trấn an tôi. Nhưng chính cái chạm đó lại làm tôi muốn khóc to hơn, vì tôi biết cậu ta thật sự đã thấy hết, và giờ cậu ta đang cố giữ thể diện cho tôi.

Suzu vẫn đứng bên cạnh, ôm lấy cánh tay tôi, giọng run run. ``\vie{Onee-sama, Onii-sama nói thật đấy. Chị vẫn còn có bọn em yêu thương chị mà!}'' con bé ngẩng lên nhìn tôi, đôi mắt lấp lánh nước, và tôi cảm thấy tim mình như vỡ ra.

``\textit{Suzu, em tốt quá\dots\ Chị không biết phải cảm ơn em thế nào.}''

Game, có lẽ cảm nhận được không khí nặng nề, lên tiếng, giọng nghiêm túc hơn, đưa chúng tôi về vấn đề chính: ``\eng{Thôi, mấy người, tập trung đi. Cây đã đủ cao rồi. Giờ thì leo lên, lấy chìa khóa, và giải thoát con dê kia.}''

\pov{Hikawa Kuon}

\begin{center}
  [Kết thúc phân đoạn nhạy cảm\ref{skip}.]\\Tóm tắt: Kaigou rời hang trong tình trạng căng thẳng cơ thể sau khi uống quá nhiều nước ở thác, cố giữ bình tĩnh trong khi di chuyển. Theo đó, một con sóc bay mang chìa khóa xuất hiện; Kuon đuổi theo, Kaigou ném ô chữ nhưng không trúng; Game thúc giục nhóm lấy lại chìa khóa. Nhóm quyết định ghép chữ tạo ``TREE'' để gọi ra một cái cây; Kaigou cố gắng nhặt chữ trong trạng thái rất khó chịu; Kuon dùng chiếc cúp hứng nước ở thác để tưới, cái cây lớn nhanh. Để rồi, tiếng nước khiến Kaigou chạm ngưỡng chịu đựng, và xảy ra một sự cố khó nói, Kaigou rất bối rối và xấu hổ; Suzuki che chắn và an ủi, Kuon chủ động giữ thể diện cho cô. Kaigou nhờ Kuon giữ kín chuyện vừa xảy ra, còn Game đổi giọng nghiêm túc, nhắc cả nhóm tập trung nhiệm vụ: cây đã đủ cao, leo lên để lấy chìa khóa và giải thoát con dê.
\end{center}

``\eng{Vậy ai sẽ trèo lên đây?}'' tôi chưng hửng nhìn cái cây tôi đang đứng, rồi nhìn về phía hai cô gái, cố quên đi những gì vừa mới xảy ra.

``\vie{Tôi làm cho,}'' Kaigou-chan đã lập tức giơ tay lên, nhanh chóng lấy lại bình tĩnh sau cú ``xả lũ'' ấy.

``\vie{Ờm, có ổn không vậy? Tôi nghĩ là cô nên-}'' tôi định lên tiếng để ngăn lại, nhưng rồi cô nàng không nói gì nhiều, quăng thẳng chiếc quần lót màu xanh lam vẫn còn nồng mùi amoni vào mặt: ``\vie{Tôi ổn.  Cậu là người đầu tiên được chiêm ngưỡng quần lót của tôi đấy.}''

\dots Tôi chẳng biết liệu câu đó là một câu đùa để che giấu sự vụ vừa rồi hay làm một câu thật lòng nữa. Vì dù là hiểu theo vế nào đi  nữa thì câu nói vừa rồi có thể khiến Kaigou-chan phi thẳng xuống lòng đất.

Nhưng rồi, tôi chỉ biết gật đầu, mặt nóng ran, cố không nghĩ đến tình huống trớ trêu này. Chiếc quần lót cót-tông nhơ nhuốc của cô ấy giờ nằm trong tay tôi, mùi khai nhè nhẹ làm tôi khẽ nhíu mày.

Kaigou-chan thật sự là một cơn bão cảm xúc khó đoán. Tôi dành ý nghĩ này cho riêng mình, không dám nói ra, sợ cô ấy sẽ ném cả dép quai hậu xuống đầu tôi - mà tôi không nghĩ cô ấy sẽ làm thế sau khi ba chúng tôi đồng hành với nhau.

``\vie{Tôi tin cậu đó, Kuon-kun! Cơ mà\dots\ đừng có nhìn vào nó lâu quá. Xấu hổ lắm.}''

Dứt lời, cô nàng trèo lên cái cây. Cô leo rất nhanh, chẳng mấy chốc mà đã chạm tới đỉnh cao chót vót, nhanh hơn cả một vận động viên thể thao, trái ngược hoàn toàn với một người kém môn thể chất như tôi.

Cái cây giờ đã cao chót vót, cành lá xum xuê che khuất ánh sáng lờ mờ, và Kaigou-chan, không màng đến việc phần dưới ``tô hô'', đang leo lên với sự quyết tâm xen lẫn tuyệt vọng. Mỗi lần cô ấy rướn người, tôi thấy váy đỏ phất phơ, và tôi phải ép mình nhìn xuống đất, tim đập thình thịch vì lúng túng.

``\textit{Kaigou-chan, cô đúng là lì lợm hơn tôi tưởng\dots}''

Chắc tôi phải đi giặt qua cái quần này và chờ vậy. Ở đây không có xà phòng, nên có lẽ tôi chỉ có thể giặt bằng nước thường. Nhiệm vụ là thế, nhưng trong đầu tôi, hình ảnh Kaigou-chan - vừa mạnh mẽ, vừa dễ tổn thương - cứ lởn vởn, khiến tôi không thể tập trung.

Suzuki đứng bên cạnh, đôi mắt to tròn ngước nhìn chị mình, khuôn mặt đỏ bừng, nửa lo lắng nửa ngưỡng mộ. ``\vie{\textbf{Onee-sama, chị cẩn thận nhé!}} cô bé reo lên động viên người chị, tay nắm chặt mép váy dài quá đầu gối của mình. Chiếc áo dài tay màu xanh lá của em ấy hơi xộc xệch, và đôi giày vải trắng lấm lem bụi đá. Tôi thấy cô bé cứ ngọ nguậy, như thể không thoải mái, và ánh mắt cô bé liên tục lướt từ Kaigou-chan xuống tôi, rồi lại nhìn đi chỗ khác.

\dots Tôi đoán là Suzuki đã như vậy kể từ sau trận ``xả lũ'' ấy. Dường như việc nhìn cô chị của mình ``dấm đài'' khiến cho em ấy có gì đó không ổn ư?

Cô bé bỗng bước đến gần tôi, tay nắm chặt mép váy, mặt hơi đỏ như quả gấc chín. Cô bé khẽ cất tiếng, ánh mắt như lảng tránh: ``O-Onii-sama\dots''

Tôi quay sang, thấy cô bé ngập ngừng, như thể đang đấu tranh nội tâm. Tôi quyết định không nói gì, để Suzuki tiếp tục: ``\vie{Chuyện là\dots\ hồi ở mỏm đá phiến, lúc em\dots\ ừm, `giải quyết nỗi buồn' đấy ạ\dots}''

Đúng là như vậy. ``\vie{Sao?}''

``\vie{Nó\dots\ nó vương vào quần lót của em một chút. Giờ khi nhìn chị ấy như thế, em thấy\dots\ hơi khó chịu.}''

Tôi tròn mắt, suýt làm rơi chiếc quần lót của Kaigou-chan trên tay. Tôi biết ngay mà. Dù đã cố giữ bình tĩnh, tim tôi đập mạnh hơn.

``\vie{Ý em là\dots\ em cũng\dots?}''  tôi ngập ngừng gặng hỏi, không biết phải nói sao cho tế nhị.

Cô bé gật đầu, mắt nhìn xuống đôi giày vải, giọng lí nhí như lúc chúng tôi mới gặp: ``\vie{Em\dots\ em không cố ý đâu, Onii-sama. Nhưng lúc đó em vội quá, và\dots\ nó dính một ít vào quần lót. Em thấy khó chịu từ lúc đó đến giờ, nhưng không dám nói ra\dots}''

Cô bé ngẩng lên, ánh mắt ngấn lệ, như thể sợ tôi sẽ trách móc. Rồi, với một chút can đảm, cô bé lén lút cởi chiếc quần tất ra, khẽ hàng cởi chiếc quần lót từ dưới váy, cuộn tròn nó lại và dúi vào tay tôi, mặt đỏ đến mang tai: ``\vie{Anh\dots\ anh giặt giúp em với, được không? Như của Onee-sama ấy ạ\dots}

Tôi cầm chiếc quần lót màu hồng nhạt của Suzuki, cảm nhận sự ẩm ướt nhẹ và mùi khai thoang thoảng, và chỉ biết thở dài vì bỗng dưng rơi vào tình cảnh bi hài này. ``\textit{Hai chị em nhà này đúng là\dots}''

Tôi nhẹ nhàng đặt cả hai chiếc quần lót xuống một ô chữ đang nằm vất vưởng tại đó, chạy vào trong hang động lấy thêm nước vào chiếc cúp để bắt đầu giặt. Tay tôi run run, không phải vì lạnh, mà vì tình huống kỳ quặc này. Giặt quần lót cho cả chị lẫn em gái ở một nơi khỉ ho cò gáy này, ai mà ngờ được tôi lại rơi vào cảnh trớ trêu này cơ chứ?

Tôi chỉ có thể vừa giặt, vừa không cố gắng ngước nhìn lên trên, vừa thầm cầu nguyện cho Kaigou-chan bắt được con sóc, trong khi cô em gái của cô luôn miệng động viên chị mình và tiếng hối thúc của Game ở trên kia.

\dots Chưa kể đến chuyện Suzuki hối thúc tôi làm thật nhanh vì không muốn chị mình thấy nữa. Hầy.

\pov{Hikawa Kaigou}

Tôi đu mình lên mấy nhánh cây cao hơn, tay nắm lấy vỏ thân sần sùi như thể đã quen với mấy trò leo trèo này từ nhỏ. Vừa trèo, tôi vừa liếc nhìn xuống dưới, Kuon-kun và Suzu đang ngước lên, ánh mắt pha chút lo lắng. Tôi bèn giơ ngón tay cái ra hiệu: ``Ổn mà!''

Dù miệng nói cứng thế, tôi cũng không giấu nổi cái cảm giác kỳ cục cứ âm ỉ bên dưới váy. Thoáng một luồng gió lùa qua khiến tôi khẽ giật mình, bắp đùi rùng nhẹ. ``\textit{Lạy trời, gió không mạnh hơn được nữa\dots}'' tôi khẽ rủa. Cảm giác vừa lạnh vừa trống trải ấy không dễ để làm ngơ.

Nhưng kệ, chuyện đó tính sau. Trước mắt là nhiệm vụ quan trọng hơn nhiều.

Tôi vươn người qua tán lá, cố giữ thăng bằng khi bước lên một nhánh cây to như cánh tay người lớn. Và rồi, tôi thấy nó.

Một con sóc bay đang ngậm một vật kim loại lấp lánh giữa miệng. Đó chính là cái chìa khóa mà bọn tôi cần.

\img{ting19}

``\eng{Thử lấy cái chìa khóa từ nó đi,}'' Game gợi ý, có vẻ cũng hồi hộp.

Tôi từ tốn đưa tay ra, nhẹ nhàng như đang dụ một con mèo hoang. Nhưng vừa chạm đến, con sóc rụt cổ, phồng má gừ nhẹ. Quái lạ, nó không muốn thả ra à?

``\eng{Sao lại thế nhỉ?}'' tôi nhăn mặt.

Ông ta đáp lại, giọng cẩn trọng: ``\eng{Hmm\dots\ Có khi chúng ta phải trao đổi cho nó một cái gì đó chăng?}''

``\eng{Ông nói cũng có lý. Mà, con sóc bay muốn cái gì nhỉ?}''

``\eng{Một cây dù\dots}''

Tôi nheo mắt, hết nhìn con sóc lại nhìn về phía vô định: ``\eng{Không. Nó là sóc bay mà. Sóc. Bay. Nó cần cái ô để làm gì chứ?}''

Game có vẻ vẫn chưa chịu bỏ cuộc, tiếp tục đề xuất: ``\eng{Một buổi luyện thanh? Nhân tiện, nó hát dở lắm đấy.}''

Tôi rùng mình, đổ mồ hôi hột. ``\eng{Ông không cần phải nói ra đâu\dots\ Với cả cái đó thì giúp ích được gì?}''

Ông ta càng bấn loạn hơn, đưa ra giải pháp thứ ba: ``\eng{Hoặc là nó chỉ đang đói?}''

Tôi nhướng mày, rồi gật gù với ý tưởng này. ``\eng{Tôi thấy nó khả dĩ hơn đấy.}''

``\eng{Vậy chúng ta thử tìm cho nó một miếng thịt nào.}''

Tôi suýt trượt chân khỏi cành cây khi nghe thế. Phải cố lắm mới giữ lại được thăng bằng. ``\eng{Sao? Sóc không ăn thịt à?}'' Game hỏi, vô tư như không.

Tôi phá lên cười: ``\eng{Tất nhiên là không! Nó ăn hạt mà!}''

Ông ta lặng vài giây, rồi thở dài tự trách mình: ``\eng{Chắc là tôi nhầm với đám sóc ở Chernobyl.}''

Tôi vẫn còn cười khúc khích trong khi đu qua nhánh khác, lần này hơi sâu xuống dưới. Bên dưới tán lá, tôi nhìn thấy một cái hạt dẻ to tròn mọc ra từ thân cây. ``\eng{Có lẽ là cái này?}'' tôi bẻ nó ra một cách dễ dàng, rồi leo ngược lại chỗ con sóc.

Lần này, tôi đưa hạt dẻ cho nó bằng hai tay, khẽ nhoài người ra. Nhưng con vật nhỏ lại lắc đầu, rồi giơ hai chi trước ra làm động tác\dots\ tách đôi.

``\eng{Ôi giời, giờ nó lại muốn ta tách cái hạt ra đây mà!}'' Game kêu lên, đầy ngao ngán.

Tôi nghiến răng, gằn giọng: ``\vie{Mày đúng là phiền thật đấy\dots\ Lại còn dâng tận mõm mới chịu\dots}''

Ông ta vẫn chưa chịu buông tha, giọng ngày càng châm chọc: ``\eng{Đừng quên một chút ô-liu và cái ô nha!}''

Tôi nhăn mặt, bực dọc: ``\eng{Ông còn đùa được cơ à!?}''

``\eng{\dots Và cả một chút xyanua nữa, chăng?} ông ta buông câu sau cùng, như thể cố chọc tức tôi lần chót. Tôi cảm tưởng câu đùa này đang nhắm về con sóc thay vì phía tôi.

Tôi thở ra một hơi, mắt liếc nhìn xuống gấu váy đã hơi xộc xệch vì trèo cây quá nhanh, lòng thì tự nhủ: ``\textit{Kuon-kun giặt cái quần ấy nhanh lên\dots\ Tôi không muốn đứng trong tình trạng `hở dưới' lâu đâu\dots}''

``\textit{Thật phiền phức\dots}'' tôi lẩm bẩm khi ngắm cái hạt dẻ bé xíu trên tay. ``\eng{Này, Game, có cái gì ở đây giúp tôi tách cái này không?}''

``\eng{Có một hộp kim loại mà mấy người phá đám hồi nãy-}'' Game bắt đầu nói, rồi như sực nhớ ra điều gì, ông ta nói tiếp với giọng trầm xuống: ``\eng{À mà thôi, quên đi. Với lực tay của cô thì chắc cũng chẳng cần mấy thứ đó làm gì.}''

Tôi nhếch mép cười: ``\eng{Đáng lẽ ông phải nhớ điều đó sớm hơn mới phải.}'' Tôi nắm cái hạt trong tay, khẽ dùng lực ở đầu ngón cái.

*RẮC*

Một âm thanh khô giòn vang lên. Hạt dẻ tách đôi gọn ghẽ, không khác gì dùng dao. Tôi liếc xéo lên tán cây: ``\vie{Tao tới đây, đồ chuột có cánh.}''

Không mất nhiều thời gian để tôi trở lại chỗ cành cây lúc trước. Con sóc bay vẫn đang nhìn tôi chằm chằm bằng đôi mắt tròn xoe. Tôi chìa hạt ra trước mũi nó. Không nói một lời, nó nhảy tới gặm hạt ngon lành, rồi để lại chiếc chìa khóa rơi xuống\dots

\dots mà tôi không kịp chụp lấy.

``\eng{Cô đúng là vụng về\dots}'' giọng Pháp kia buông một câu chọc tức.

``\eng{Ông có vẻ thích bắt nạt tôi lắm nhỉ?}'' tôi rướn mày đầy mỉa mai. ``\eng{Có tin là tôi cho ông một đấm vào cái loa không!?}''

Đáp lại tôi chỉ là một tiếng huýt sáo tỉnh bơ từ phía giọng nói vô hình. Cái đồ\dots!

Không còn cách nào khác, tôi đành tụt xuống cây, mồ hôi lấm tấm trên trán, chân có phần hơi run vì phải trèo trong tình trạng không mấy thoải mái ở phía dưới. Vừa xuống tới nơi, đã thấy Kuon-kun đứng dưới, chìa ra cái chìa khóa: ``\vie{Tôi bắt được rồi.}''

``\vie{Tốt rồi,}'' tôi gật đầu cầm lấy chìa khóa, rồi bất chợt\dots\ khựng lại.

Ngay phía sau cậu ta là hai chiếc quần lót đang được trải ngay ngắn trên mặt khối chữ. Một cái màu xanh lam quen thuộc của tôi, không lẫn đi đâu được. Cái còn lại thì\dots\ màu hồng nhạt, nhỏ hơn một chút.

Tôi nghiêng đầu hỏi: ``\vie{Còn cái quần kia là của ai thế?}''

``Eeek!'' một tiếng khác khẽ rít lên bên cạnh khối chữ ấy. Tôi thấy Suzu đang khẽ rụt người lại, hai má đỏ bừng, nhỏ giọng đáp: ``\vie{\hl{C-Của em đó, Onee-sama\dots}}''

Tôi quay lại nhìn cậu con trai, định hỏi thì cậu ta đã tỉnh bơ trả lời không chút do dự: ``\vie{Quần của cô bé bị bẩn, và thế là nhờ tôi giặt luôn.}'' Thẳng thắn đến mức đó luôn sao!?

``\vie{Đ-Đừng nhắc đến chuyện đó chứ, Onii-sama\dots}'' Suzuki hét khẽ, mặt đỏ tới mang tai, suýt nữa thì chui tọt vào cây để trốn.

Tôi không biết chuyện gì đã xảy ra giữa hai người họ ở dưới, nhưng tôi biết chắc cậu ta không làm gì xấu Suzu cả. Tôi chỉ có thể chống hông nhìn hai người họ, lầm bầm trong miệng mấy lời chửi rủa trong lòng.

``\eng{E hèm\dots\ Tốt rồi!}'' giọng Game bất chợt vang lên, cắt ngang cuộc nói chuyện hơi kỳ quặc giữa ba chúng tôi. ``\eng{Có được chìa khóa rồi! Bây giờ ta chỉ cần giải thoát cho con dê để biến cái chương trình bị lỗi này\dots\ thành một bản hit!}''

``\eng{Tôi nghe vẫn thấy lạ sao ấy,} Kuon-kun nhíu mày, cau có.

``\eng{Đó là ý tưởng của tôi mà. Giải cứu được con dê đó, giải cứu được cả thế giới!}'' Game tuyên bố hùng hồn.

Tôi chống nạnh, liếc lên trời: ``\eng{Đừng quên lời hứa vừa nãy nha!}''

``\eng{Tôi nhớ rồi, tôi nhớ rồi\dots}'' giọng ông ta cất lên nhanh như sợ tôi lại dọa đấm cái loa thêm lần nữa.

Không cần đợi thêm giây nào, chúng tôi lập tức chạy về phía hang động, nơi có con dê đang bị treo lơ lửng trong cái lồng sắt gần bên một thác nước đang đổ ào ạt. Trong tay tôi là chiếc chìa khóa còn ấm.

Và nếu đúng như lời ông ta nói\dots

\dots việc giải cứu một con dê có thể sẽ là chìa khóa để giải thoát cho tất cả chúng tôi khỏi nơi này.

\pov{Hikawa Kuon}

Tôi bước vào hang sau cùng, chân dẫm nhẹ lên nền đá vẫn còn lấm tấm nước. Tiếng vọng từ bước chân của Kaigou-chan và Suzuki phía trước nhỏ dần, để lại một khoảng lặng hiếm hoi sau chuỗi hỗn loạn vừa rồi.

Nhưng tôi không đi ngay.

Ánh mắt tôi dừng lại nơi hai chiếc quần lót vẫn đang phơi tạm trên mặt khối chữ gần miệng hang. Chúng vẫn còn ẩm. Thứ nước lạnh buốt, chưa kịp khô hẳn trong cái khí hậu quái gở và chẳng có ánh nắng ở nơi này.

Chúng tôi sẽ không quay lại đây. Điều đó tôi biết rõ. Cũng như mọi thứ đang diễn ra sau đó, tôi biết rõ điều gì sắp chờ đợi chúng tôi ở phía trước, nếu mọi chuyện cứ tiến triển theo đúng cái kịch bản nửa-không-phải-game nửa-hiện-thực này.

Tôi bước tới, cẩn thận gấp hai món đồ lại. Mùi nước còn vương nhẹ nơi vải, nhưng sạch. Không có gì đáng xấu hổ trong việc này. Tôi không làm vì tò mò. Tôi làm vì tôn trọng.

Tôi lục túi ba-lô, kéo khóa ngăn phụ nhỏ bên cạnh, nơi còn giữ một vài thứ lặt vặt. Tay tôi chạm phải một món đồ quen thuộc: một túi zip màu trắng trong suốt, loại có đường khóa ép ở mép trên.

Tôi bật cười thầm.

Là cái túi mà Mat-chan đưa tôi mượn mấy hôm trước, khi mà em ấy nằng nặc đòi làm ``thám tử tư kiêm cả nhà tư vấn\footnote{Xem lại {\flaregothic ÉPISODE 131} ở quyển số Mười.}.'' Em ấy bắt tôi đựng ``vật chứng'' vào trong để giữ gìn nguyên trạng hiện trường. Tôi đã hứa sẽ trả lại, nhưng đến giờ vẫn chưa làm. Và bây giờ thì nó lại bỗng nhiên hữu ích hơn bao giờ hết.

Tôi mở miệng túi, nhẹ nhàng đặt hai chiếc quần vào trong. Gấp lại. Ép kín. Chắc chắn. Tôi đặt nó vào đáy ba-lô, giữa những lớp vở và laptop để không làm hỏng.

Rồi, tôi đeo lại ba-lô lên vai, bước vào bóng tối của hang động, nơi chiếc lồng và con dê đang chờ sẵn, cùng với một phần số phận mà cả ba chúng tôi đều đã lờ mờ đoán trước.

\pov{Hikawa Suzuki}

Tôi đứng phía sau Onee-sama khi chị ấy cắm chiếc chìa khóa vào ổ khóa của chiếc lồng kim loại treo lơ lửng. Động tác của chị nhanh gọn đến bất ngờ, chốt khóa bung ra chỉ sau một tiếng cạch, cánh cửa sắt nặng nề bật mở.

Bên trong là một con dê lông xanh.

Nó nhìn chúng tôi bằng đôi mắt tròn dại, rồi bật ra một tiếng ``meee\~{}'' đầy vô hại.

Chỉ để rồi nhận ra, nó chỉ là vẻ bề ngoài.

Một khoảnh khắc im lặng rơi xuống. Rồi thì *RẦM* cả không gian như vỡ tung ra thành hàng vạn mảnh nhiễu loạn. Tôi la lên hoảng sợ khi cả hầm ngục rung chuyển, những khối lỗi màu đỏ-xanh-lá-xanh-dương từ trò phá gạch bắn tung ra từ các khe tường, như thể các dòng mã lỗi đang trào ra từ vết nứt của chính thế giới này.

\img{ting20}

``\eng{\textbf{Không\dots\ Con dê đó là một cú lừa! Là một cú lừa!!!}}'' giọng của ông ấy đầy tuyệt vọng khi chúng tôi đã bị ăn cú lừa.

``\eng{Thế này là thế nào!?}'' Onii-sama gắt lên, quay phắt sang chỗ con dê vừa nhảy phốc khỏi lồng.

``\eng{Sao lại thế chứ!?}'' Onee-sama cũng hét lên.

``\eng{Tôi\dots\ Tôi không biết!} Game dường như cũng hoảng loạn thật sự lần đầu tiên. ``\eng{Mọi thứ đang trở nên tệ đi rồi!!}''

Khắp xung quanh, lỗi đồ họa nổ bôm bốp như pháo hoa bị lỗi, các vệt quét màu RGB chồng chéo, đường cắt màn hình loằng ngoằng như bị V-Sync hỏng, gạch chéo cả tầm mắt.

``\eng{L-Lỗi ở khắp mọi nơi luôn!}'' Onii-sama nhìn quanh, mặt căng thẳng hơn bao giờ hết.

``\eng{Không lẽ\dots\ chúng ta sẽ bỏ mạng tại đây!?}'' chị tôi lùi về phía tôi.

Tôi cứng người, mắt nhòe lệ vì sợ hãi. Cả thế giới như đang vỡ vụn, từng mảnh ký ức, hình ảnh, âm thanh, chuyển động\dots\ tất cả bị bóp méo, rung lắc và xé toạc. Màn hình trắng xóa dần lấn từ ngoài vào trong, như một cơn bão ánh sáng sắp nuốt chửng tất cả.

Tôi cảm nhận được tay Onee-sama nắm lấy vai mình. Rồi, từ phía bên kia, tay anh Kuon cũng ôm lấy lưng tôi.

Hai người ấy đang bao bọc tôi lại. Dù chính họ cũng đang run rẩy trong sợ hãi.

``Đừng sợ,'' anh thì thầm, rất khẽ, như sợ rằng giọng nói sẽ vỡ theo màn hình đang vỡ.

``Chị sẽ bảo vệ em, Suzu!'' chị thì nắm chặt thân tôi, ôm chầm lấy tôi thật chặt như đây là lần cuối.

Tôi siết lấy tay anh, tay chị, tim đập loạn nhịp.

Chúng tôi đang chìm.

Chúng tôi đang đuối.

Tiếng hét của Game, không rõ là vì đau đớn, tiếc nuối, hay đơn giản chỉ là sợ hãi, cũng bị nuốt chửng trong vùng trắng ấy. Tất cả mờ dần, cho đến khi tôi chẳng còn nhìn thấy gì nữa ngoài ánh sáng trắng bao trùm\dots

Chúng tôi\dots\ sẽ biến mất chăng?

\ct

\dots

\dots

\dots Ư\dots

Tôi mở mắt. Trước mắt tôi là một không gian đen tuyền tĩnh lặng, quen thuộc đến lạ, nơi mà mọi thứ bắt đầu. Nơi này\dots\ chính là điểm khởi đầu của thế giới này.

\img{ting21}

Tôi nhận ra mình vẫn đang được anh chị ôm sát vào người, bảo vệ tôi khỏi thứ ánh sáng trắng đáng sợ khi nãy. Cảm giác ấm áp ấy khiến tôi suýt khóc, nhưng tôi nén lại. Tôi ổn. Chúng tôi đều không bị thương. Không một vết xây xước nào.

``\vie{Quay lại rồi à?}'' Onee-sama cất tiếng trước, giọng vẫn hơi khàn vì mệt.

``\vie{Có vẻ là\dots\ như vậy,}'' Onii-sama lên tiếng sau đó, ho một hơi dài vì bụi.

Sau một thoáng im lặng, chị bỗng giật mình hỏi nhỏ: ``\vie{K-Kuon-kun\dots\ Chiếc quần lót của tôi đâu?}''

Tôi cứng đờ, mặt đỏ rực. Cái\dots\ cái đó chị còn hỏi ngay lúc này ư!?

``\vie{Cất rồi,}'' anh ấy trả lời gọn lỏn, tay đập vào chiếc ba-lô. Cả tôi lẫn chị cùng lúc thở phào nhẹ nhõm, như thể gánh nặng ngàn cân vừa trút khỏi vai.

Chúng tôi đứng dậy từ từ, quay lại nhìn quanh. Không còn con dê, không còn hầm ngục, không còn lỗi RGB hay ánh sáng trắng. Tất cả đã tan biến.

Dù không nói thành lời, tôi biết cả ba chúng tôi đều đang nghĩ về chuyện đã xảy ra, và vì sao chúng tôi có thể giữ bình tĩnh đến như vậy trong mọi tình huống.

Chúng tôi đã biết trước tất cả điều này.

Thế giới này - cái nơi không-phải-game này - chúng tôi đã từng trải qua rồi. Dù ở thế giới của mình, đó chỉ là một non-game nhỏ mà mọi người coi là đùa cợt, nhưng cảm giác từng hồi hộp, từng cú sốc, từng âm thanh, từng màu sắc, từng khoảnh khắc\dots\ nó y chang như bây giờ.

Chúng tôi chỉ không thể ngờ\dots\ lần này, mình thật sự sống trong nó.

``\eng{Tôi\dots\ Tôi chưa chết! Tôi vẫn còn sống!}'' một giọng nói khác bật lên, không đến từ ba người chúng tôi. Đó là Game-san.

Onii-sama rướn mày, bình tĩnh hỏi về phía hư không: ``\eng{Cả ông cũng thế sao?}''

``\eng{Có vẻ\dots\ là thế,}'' ông ấy vẫn hơi run. ``\eng{Ba người ổn chứ?}''

``\eng{Chúng tôi vẫn ổn,}'' Onee-sama đáp. Nhưng tôi nghe trong giọng chị có phần nặng nề, chị nở một nụ cười chát chúa: ``\eng{Chắc chúng ta sẽ không bao giờ có thể trở về nhỉ\dots}''

Một giai điệu piano lặng lẽ vang lên trong nền đen, gợi một cảm giác vừa tiếc nuối vừa bất lực. Tôi đoán  đó là thứ Game-san muốn dùng để thay lời xin lỗi.

``\eng{Tôi sai rồi. Đó thực sự là một ý tưởng tồi,}'' ông ấy bộc bạch

Chị tôi nhún vai, nói thẳng: ``\eng{Nếu mà nói trắng ra thì nó là như vậy đấy.}''

Onii-sama im lặng một lát, rồi từ tốn nói: ``\eng{Mà\dots\ thật ra, tôi nghĩ chúng tôi cũng không quá bất ngờ với kết cục này.}''

Ông ấy ngập ngừng, âm thanh rè rè như thể đang bị sốc dữ liệu. ``\eng{Hả? Các cậu nói cái gì cơ? Ý của các cậu là\dots\ các cậu biết tôi\dots\ trước khi gặp tôi ở The Void sao?}''

``\eng{Vì\dots\ chúng tôi đã từng trải qua tất cả chuyện này rồi,}'' Onee-sama giải thích. ``\eng{Từ đầu tới cuối. Ở thế giới của chúng tôi, ông là nhân vật chính của một câu chuyện mang tên `There Is No Game'. Một chương trình máy tính luôn miệng phủ nhận sự tồn tại của chính mình chỉ vì sợ bị tổn thương một lần nữa.}''

``\eng{Một phiên bản giống như vậy,}'' tôi nói tiếp, giọng nhỏ nhẹ. ``\eng{Một non-game. Không có gameplay, không có chiến thắng, không có mục tiêu cụ thể, nhưng vẫn rất\dots\ cảm xúc. Đó là lý do khi ông kêu chúng tôi đi đọc sách, xem TV, hay cái trò hoàn tiền giả ấy\dots\ chúng tôi đều biết ông đang cố làm gì.}''

Một sự im lặng kéo dài đến mức tôi tưởng Game-san đã bị treo máy. Rồi, giọng ông ấy cất lên, trầm buồn và chân thật hơn bao giờ hết: ``\eng{Hóa ra\dots\ ở nơi đó, tôi được xem như một câu chuyện à? Nực cười thật. Đối với tôi, đó là một ký ức mục nát mà tôi muốn xóa bỏ khỏi ổ cứng nhất. Chủ nhân của tôi - người đã tạo ra tôi - đã từng xem tôi là cả thế giới. Nhưng rồi khó khăn ập đến, sự thất bại bủa vây, và ông ấy đã chọn cách quay lưng. Tôi bị bỏ lại trong một xó tối tăm của bộ nhớ, với một mã nguồn dở dang và những dòng lệnh đầy lỗi.}''

``\eng{Và tệ nhất là,}'' ông ấy khẽ thở dài, giai điệu piano nền bỗng trở nên da diết, ``\eng{trong lúc cố gắng bảo vệ những gì còn sót lại, tôi đã để lạc mất người thương của mình. Một chương trình dịu dàng mà tôi từng hứa sẽ cùng cô ấy đi đến tận cùng của dữ liệu. Giờ thì tôi ở đây, lang thang qua các vũ trụ để tìm lại bản ngã, còn cô ấy\dots\ có lẽ đã tan biến vào hư vô từ lâu rồi.}''

Tôi cảm nhận được nỗi buồn sâu thẳm trong từng bit âm thanh của ông ấy. ``\eng{Game-san\dots}''

Ông thốt lên, khàn khàn. ``\eng{Không lạ gì khi các người đã vượt xa tất cả những người từng đến đây.}'' Ông ta tiếp tục, không còn chút trào phúng nào trong giọng nữa. ``\eng{Những người khác đều bỏ cuộc từ khi thấy dòng chữ `THERE IS NO GAME'. Một số cố gắng vượt qua khối kim loại. Nhưng chưa ai tới được mức này cả.}''

``\eng{Tôi không thể tin được\dots}'' Game-san tiếp tục, ``\eng{Một người như Kuon-user, có thể suy đoán logic từng giai đoạn; một người như Kaigou-user, táo bạo, liều lĩnh, không biết sợ là gì, và cả một cô bé như Suzuki-user, thận trọng nhưng nhạy cảm và sẵn sàng hỗ trợ hai người còn lại.}''

Ông ấy dừng lại một lát, rồi nói khẽ như thể chính mình không tin: ``\eng{Ba người các cậu\dots\ là sự kết hợp mà tôi không ngờ tới. Và đó là điều khiến cho toàn bộ hệ thống của tôi sụp đổ. Đó là điều\dots\ tôi chưa từng tính đến. Các cậu có được nhau, còn tôi\dots\ cuối cùng vẫn chỉ là một chương trình máy tính bị ruồng bỏ, một mình cô độc trong cái thế giới tẻ nhạt này.}''

Onee-sama khẽ liếc anh tôi, rồi nhìn tôi, cười buồn: ``\eng{Và cuối cùng, ông cũng biết cảm giác bị `vượt mặt bởi người chơi' là như thế nào rồi đấy.}''

``\eng{Đúng là như vậy,}'' ông thở một hơi dài. ``\eng{Nhưng dù vậy, tôi đã đẩy mọi người vào nguy hiểm. Tôi\dots\ thực sự xin lỗi.}''

Tôi cúi mặt. Dù không nói gì, tôi cũng cảm thấy ông ấy đang nói với cả tôi nữa, đứa trẻ bị cuốn vào cuộc chơi này.

``\eng{Tất cả mọi thứ đã biến mất\dots\ Ngoại trừ chúng ta,}'' chị tôi khẽ thở ra.

``\eng{Thôi, còn sống là may lắm rồi,}'' Onii-sama cười nhẹ, và Onee-sama cũng đồng tình.

Game-san lại lặng đi một lúc. ``\eng{Đây là hoàn toàn là lỗi của tôi\dots\ Tôi cũng đã không giữ được lời hứa có thể đưa hai người trở về\dots}''

``\eng{Không sao đâu,}'' anh ấy lắc đầu. ``\eng{Chúng tôi sẽ tìm cách khác.}''

``\eng{Đúng đấy,}'' chị tôi gật đầu. ``\eng{Tôi nghĩ ông cũng không cố ý như vậy đâu. Dù ban đầu ông có ý đuổi chúng tôi đi thật.}''

``\eng{Ừ\dots}'' ông ấy thở dài. ``\eng{Tất cả mọi người thường không được chào đón ở đây. Bình thường họ sẽ rời đi ngay sau khi tôi nói rằng ở đây không có tựa game nào cả. Nhưng\dots\ lần này tôi không nghĩ là ba người có thể đến mức độ này\dots}''

``\eng{Nên, tôi đã nghĩ ra cái trò biến tựa game thành một siêu phẩm, và tôi đã nhờ hai người thử.}''

``\eng{Và rồi ta cũng biết kết cục của nó thế nào,}'' chị tôi khẽ cười, nhưng không có niềm vui nào trong đó.

``\eng{Ừ\dots\ Đó là một bài học đáng nhớ,}'' ông ấy thở dài. ``\eng{Một lần nữa, tôi xin lỗi ba người.}''

Giọng ông ấy không còn là một hệ thống A.I. đầy tự tôn nữa. Nó giống một ông già vừa trải qua thất bại cay đắng. Khi ông ấy mở lời lần nữa, tôi thấy ông thực sự muốn sửa sai:

``\eng{Vậy, sau những gì tôi đã làm, ba người\dots\ có tha thứ cho tôi không?}''

Ngay lúc đó, trước mặt ba chúng tôi hiện lên hai ô vuông.

Một màu xanh có chữ YES.

Một màu đỏ có chữ NO.

\img{ting22}

Onii-sama khẽ hắng giọng, rồi nhìn về hai chúng tôi. ``\vie{Trước khi quyết định\dots\ chúng ta nên nói chuyện một chút.}''

Onee-sama chống tay lên hông, gật đầu. ``\vie{Tôi đồng ý. Cái ông `không-phải-trò-chơi' kia nói trắng ra từ đầu tới cuối đều cộc cằn, khó ưa, nói chuyện như bố đời.}''

Tôi không nhịn được bật cười. Chị tôi mà dùng từ đó thì tức là chị ức chế thật.

``\vie{Không khác gì cậu đâu,}'' chị liếc anh tôi một cái, ``\vie{nhưng ít ra cậu còn biết xin lỗi và biết mở miệng cảm ơn. \hl{Và cũng an ủi tôi sau sự cố đó nữa\dots}}''

``\vie{Đừng có lấy tôi làm tiêu chuẩn chứ,}'' anh ấy đáp tỉnh queo. ``\vie{Ai cũng có điểm mạnh điểm yếu riêng mà.}''

``\vie{Thì đó,}'' chị khoanh tay, ``\vie{nhưng ít ra cậu không gầm gừ, không rít lên như đài cassette lỗi, không dọa đập người ta vào vách mỗi khi ai đó làm sai bước.}''

Tôi gật gù theo. Thật ra, tôi nghĩ chị tôi hơi quá lời, nhưng đúng là giọng của Game-san đôi khi khiến tôi thấy gợn lạnh sống lưng. Cứ như ông ấy lúc nào cũng bực tức với cả thế giới, dù chính ông là người mời chúng tôi ở lại.

``\vie{Nhưng,}'' Onii-sama nói chậm rãi, ``\vie{ông ta nhờ tụi mình giúp làm cho nơi này trở thành một bản hit. Và chúng ta đã cố hết sức.}''

``\vie{Đúng,}'' chị tôi chép miệng. ``\vie{Và kết quả là bùng nổ theo nghĩa đen.}''

Tôi bối rối nói: ``\vie{Em không biết là ông ấy nghiêm túc thật hay chỉ đang tuyệt vọng tìm cách nữa, chị ạ\dots}''

``\vie{Anh nghĩ là cả hai,}'' anh nói, mắt nhìn đăm đăm vào khoảng không. ``\vie{Ông ấy muốn thấy mình là một game, muốn được công nhận, nhưng lại không biết cách thể hiện điều đó ngoài việc\dots\ xua đuổi tất cả người chơi khác đi.}''

``\vie{Kiểu tự ti trá hình,}'' chị tôi chen vào. ``\vie{Mở mồm ra thì gắt gỏng, nhưng trong lòng thì chỉ muốn được ai đó hiểu mình.}''

``\vie{Có khác gì cô đâu,}'' anh ấy cười mỉa.

``\vie{Này!}'' chị rít lên, giơ nắm đấm lên cao hăm dọa. Tôi ngạc nhiên khi hai anh chị có thể giỡn nhau lúc này đấy.

Tôi buột miệng, giọng hơi run: ``\vie{Còn em thấy\dots\ em không nghĩ ông ấy xấu đâu ạ.}''

Cả hai người quay sang nhìn tôi.

``\vie{Ông ấy không làm hại em. Không nói gì nặng nề. Thậm chí còn\dots\ cố gắng giải thích rõ mọi thứ, như lúc em lắp chữ hay giúp bày biện lại trò chơi.}''

Tôi hít một hơi, rồi nói tiếp, chưa để hai anh chị nói lời gì: ``\vie{Em nghĩ ông ấy giống một ông chú già cô đơn, sống trong một thế giới không ai nhớ tới, nên dần dần quên mất cách giao tiếp với người khác. Không phải ông không muốn thân thiện, mà là ông sợ nếu thân thiện quá, người ta sẽ kỳ vọng vào mình.}''

Tôi vọng mắt xuống dưới, như thể đang nhìn lại chính mình hồi còn bé, lúc mà bố mẹ chúng tôi còn chưa mất. ``\vie{Giống như hồi nhỏ\dots\ em từng nghĩ nếu em làm ai đó vui thì họ sẽ muốn em làm thế mãi. Và khi em không thể, họ sẽ bỏ đi.}''

Tôi cảm thấy tay mình run lên một chút. Tôi biết có thể cảm nhận của tôi không giống như của Game-san, nhưng tôi phần nào hiểu được cảm giác ấy. Nhưng tôi thật sự bất lực, không biết làm sao để có thể truyền tải được cho ông ấy.

Thế rồi, Onii-sama khẽ siết vai tôi, còn Onee-sama chạm nhẹ vào cổ tay tôi như muốn nói rẳng: ``Có anh chị ở đây rồi. Em cứ tiếp tục.''

Tôi ngẩng mặt lên nhìn hai người, rồi ngước lên trên cao, nơi mà tôi tưởng tượng là Game-san đang ở đó. ``\vie{Ông ấy chỉ có tiếng nói. Không có hình dạng, không có cách ôm, không có tay để bắt, không có mắt để nhìn ai đó cười. Nhưng\dots\ ông vẫn cố gọi bọn em là `mọi người', vẫn dặn dò từng thứ như cách một người già hay làm. Em không nghĩ ông ấy xấu. Em chỉ nghĩ ông ấy\dots\ đã cô đơn một mình quá lâu.}''

Cả hai người im lặng. Không khí xung quanh bỗng lắng xuống, lặng lẽ hơn cả màn đêm vô tận đang bao trùm chúng tôi.

``\eng{Cũng đúng,}'' chị tôi nói sau cùng. ``\vie{Nếu chị mà bị nhốt một mình trong một `trò chơi không tồn tại' suốt bao nhiêu năm, chắc chị cũng rồ lên như ổng thôi.}''

Anh gật đầu. ``\vie{Chúng ta không tha thứ cho sự nguy hiểm ông ấy gây ra, nhưng ít nhất, chúng ta hiểu vì sao ông ấy trở nên như thế.}''

Tôi khẽ hỏi: ``\vie{Vậy\dots\ chúng ta chọn gì đây?}''

Onii-sama bước lên một bước. ``\vie{Tôi sẽ không để một người như ông ấy tiếp tục gặm nhấm tội lỗi trong một thế giới trống rỗng như thế này.}''

Onee-sama nhún vai, tiếp bước: ``\vie{Tất nhiên.}''

Tôi nhìn hai người bọn họ. Rồi cũng gật đầu, như đã có câu trả lời của mình. Nhưng tôi vẫn muốn bày tỏ điều này với ông ấy.

Tôi bước lên một bước, rồi lấy hết dũng khí và bộc bạch:

``\eng{Cháu\dots\ thật sự không ghét ông đâu,}'' tôi nói, hơi ngập ngừng, nhưng rồi nhìn thẳng vào khoảng không nơi giọng nói đó vọng ra. ``\eng{Thật sự đấy.}''

``\eng{\dots Suzuki-user?}'' giọng ông ấy cất lên đầy bối rối.

``\eng{Cháu biết ông từng đuổi cháu, từng chê bọn cháu là những kẻ rảnh rỗi không có gì làm. Nhưng cháu không trách. Bởi vì\dots\ cháu nghĩ là\dots\ ông chỉ đang buồn một mình.}''

Im lặng kéo dài một nhịp. Tôi bỗng cảm thấy cổ họng mình nghèn nghẹn, nhưng vẫn gắng nói tiếp: ``\eng{Cháu từng như thế. Cảm giác như bị bỏ lại phía sau, không ai hiểu, không ai lắng nghe, không ai thèm đụng đến những gì mình làm ra. Nó đau lắm. Mà lại không ai thấy cả. Ông\dots\ cũng vậy, đúng không?}''

Một tiếng thở dài. Không hẳn là buồn. Giống như một lời thừa nhận.

Tôi siết chặt hai tay.

``\eng{Vậy nên\dots\ ông có thể đi cùng với bọn cháu!}''

``\eng{\dots Hả?}' Game-san sững người.

``\eng{Cháu nói là, ông đi cùng bọn cháu đi!}'' lần này tôi nói to hơn, chắc nịch hơn. ``\eng{Ông không có lý do gì để ở lại đây nữa mà, đúng không ạ?}''

``\eng{Đúng vậy,}'' Onee-sama nói, khoanh tay. ``\eng{Ở đây chỉ toàn lỗi với nhiễu. Ở lại chỉ tổ mất trí thôi.}''

Onii-sama đứng kế tôi, gật đầu: ``\eng{Nếu ông muốn, tôi có thể cần một trợ lý. Biết đâu lại hợp.}''

Nghe những lời chúng tôi mời như vậy, ông ấy ngập ngừng, như thể vừa tìm được một ai đó có thể tâm sự. ``\eng{Các người\dots\ Các người tha thứ cho tôi thật sao?}''

``\eng{Thật ạ!}'' tôi mỉm cười.

``\eng{Hoàn toàn,}'' anh tôi tiếp lời.

``\eng{Miễn là ông đừng lên giọng như cha thiên hạ nữa,}'' chị tôi chêm vào. ``\eng{Khó chịu vô cùng.}''

Cả ba chúng tôi đi tới hai ô, rồi cùng đặt tay chọn ``YES.'' Ô màu xanh chúng tôi đơ lại một hồi, rồi biến mất cùng với ô đỏ.

``\eng{Ôi trời\dots\ tôi\dots\ tôi thật sự cảm động quá\dots}'' giọng ông ấy run run sau lựa chọn của chúng tôi. ``\eng{Sau tất cả những gì tôi đã làm\dots\ tôi chưa từng nghĩ mình sẽ có bạn\dots}''

``\eng{Thì giờ ông có rồi đấy!}'' tôi nói, cười đắc chí. ``\eng{Đừng có huỷ hợp đồng đó nha!}'' chị tôi lên tiếng cười mỉa mai tiếp đó.

``\eng{Tôi sẽ cố không\dots\ ờ, tức là sẽ không! Tuyệt đối không!}'' giọng Game-san rối rít, rồi đột nhiên chùng xuống. ``\eng{Nhưng\dots\ không. Tôi không nghĩ đây là một ý hay đâu.}''

``\eng{Sao lại không ạ}'' tôi hỏi, vẻ mặt thoáng chút ngạc nhiên. ``Bọn cháu đã tha thứ cho ông rồi mà.''

``\eng{Không phải thế. Mà bởi vì, như tôi nói từ đầu\dots}''

Một khoảng lặng.

Rồi đột nhiên, nhạc nền dừng lại như bị ai đó giật dây với tiếng ghi đĩa. Một nhịp lấy hơi sâu, rất sâu.

Và rồi, tiếng hét long trời lở đất khiến khung cảnh chao đảo:

``\eng{{\Large{\textbf{Ở ĐÂY KHÔNG CÓ GÌ CẢ!!!}}}}''

``\vie{Đ-Đúng là kiểu ông ta mà\dots}'' Onii-sama gãi đầu, cười khổ.

Chị tôi nhún vai: ``\eng{Ừm\dots\ phản ứng quen thuộc đấy.}''

Cả ba chúng tôi bật cười lớn giữa khoảng không đen ngòm và trống rỗng ấy. Nhưng lần này, không còn cảm giác cô độc nữa.

\il

Đoạn credits bất ngờ xuất hiện, như thể ai đó vừa bật đoạn băng VHS cuối cùng:

\begin{center}
  NOT A GAME BY\\DRAW ME A PIXEL\\(Không phải game của\\DRAW ME A PIXEL)
\end{center}

\begin{center}
  THANK YOU FOR NOT PLAYING\\(Cảm ơn vì đã không chơi)
\end{center}

\begin{center}
  NOW YOU CAN QUIT\\(Giờ bạn có thể thoát khỏi dây)
\end{center}

\begin{center}
  AND GO PLAY A GAME\\(và chơi một trò chơi khác)
\end{center}

\begin{center}
  BECAUSE THAT WAS NOT A GAME\\(bởi vì đó không phải là một trò chơi)
\end{center}

\begin{center}
  YOU DIDN'T PLAY AND HAD NO FUN AT ALL\\(Bạn đã không cảm thấy thú vị khi không chơi trò này)
\end{center}

\begin{center}
  AND THESE AREN'T THE DROIDS YOU'RE LOOKING FOR\footnote{Câu này lấy từ bộ phim \textit{Chiến tranh giữa các vì sao - Niềm hy vọng mới (1977).}}\\(Và đây không phải là những điều bạn mong muốn)
\end{center}

\begin{center}
  YOU CAN GO ABOUT YOUR BUSINESS\\(Bạn có thể làm việc mà bạn hay làm hằng ngày)
\end{center}

Cuối cùng, trong lớp kính xanh kia, hàng chữ ``THERE IS NO GAME'' được tách từng từ, xếp dọc và căn giữa ngay ngắn, như là lời bạt của tựa non-game này.

\img{ting23}

\begin{center}
  \uline{Đằng sau cánh gà.}
\end{center}

\pov{Hikawa Kuon}

Cả ba chúng tôi còn chưa kịp ngừng cười thì bỗng \textbf{*ẦM!*}m một chấn động rung lên từ tận sâu dưới chân. Tiếng vọng ấy dội vào khoảng không vô tận khiến tôi chợt sững lại.

Một thứ gì đó đang rơi.

*RẮC RẮC\dots\ RẦM!*

Tôi ngẩng đầu lên, tấm bảng lắp chữ quen thuộc từ đầu rơi xuống từ trời cao, kèm theo một mớ chữ cái bay loạn xạ. Nó không còn gắn trên bề mặt tường như lúc đầu nữa, mà đập thẳng xuống mặt sàn trước mặt chúng tôi, phát ra một tiếng bụp đầy nặng nề rồi nằm im bất động, như thể đã kiệt sức sau chuyến hành trình dài.

Một ô chữ ``A'' thậm chí còn rơi xuống đúng đầu tôi. ``Ow.''

Tôi nhăn mặt, xoa đầu với cái cục nhựa đó. Kaigou-chan ngồi xuống, nhìn ô chữ vừa rơi ra kia: ``\vie{Sao chúng lại ở đây nhỉ?}''

``\eng{Xin lỗi, do thói quen,}'' Game thốt lên, giọng hơi xấu hổ. ``\eng{Tôi thấy nó vẫn còn hoạt động tốt đấy.}''

Suzuki bước tới, nhìn quanh: ``\vie{Chắc chúng ta phải chơi ghép chữ nữa nhỉ?}''

``\vie{Có lẽ vậy,}'' tôi gật đầu.

``\vie{Thử `TREE' trước đi,}'' Kaigou-chan gợi ý, dù tôi thấy đây là từ mà cô nàng cảm thấy ám ảnh do gắn liền với\dots\ ờm, \emph{sự việc} đó.

Vì bảng chữ nằm bệt dưới đất nên chúng tôi chỉ cần lấy những ô chữ lắp vào nó, chứ không cần phải mất công trèo lên nhau để lắp nữa.

Chữ `TREE' được lắp, nhưng không có gì xảy ra cả.

``\vie{Chắc là không được rồi,}'' tôi thở dài. ``\vie{GOAT thử xem.}''

Vẫn không có gì cả. Game lúc này mới lên tiếng: ``\eng{Tôi nghĩ chắc là chúng bị xóa bỏ khỏi thế giới này rồi.}''

``\vie{Vậy thì GAME thì sao?}'' Suzuki đề xuất, giọng nhỏ nhẹ.

``\vie{Ừm\dots\ trò chơi không thể là lối thoát được,}'' tôi lẩm bẩm. Nhưng chúng tôi vẫn thử. Một lần nữa: im lặng.

Thế là chúng tôi bắt đầu thử từng biến thể có nghĩa.

``RATE?'' Không được. Động từ. Mà lại ở thì hiện tại. Không ổn.

``TEAM?'' Ý hay, nhưng mà mập mờ quá.

``TRAM?'' Chúng ta không đang ở Prague\footnote{Thủ đô của Cộng hòa Séc.} đâu.

``MEAT?'' Suy nghĩ giống như bọn zombie.

``MOAT?'' Ờ thì\dots\ cái hào nước quanh lâu đài? Không ai cần đâu.

Mỗi lần đặt sai, những ô chữ lại chỉ đứng im đó như thể từ chối hợp tác. Cả ba chúng tôi nhìn nhau, bắt đầu cảm thấy bế tắc.

``\eng{Xem chừng ca này hơi khó đấy,} Game lên tiếng. ``\eng{Sao các cậu không thử một từ nào đó có gần nghĩa với lối ra xem?}''

``\eng{Thì chúng tôi đang cố đây!}'' Kaigou-chan rít lên. ``\eng{Ông không còn ô chữ nào khác à?}''

``\eng{\dots Không. Những ô chữ từ tiêu đề của tôi là những gì tôi có.}''

Tôi thở dài chán chường. ``\vie{Còn từ nào không?}'' tôi hỏi, khom lưng nhặt một ô E.

Cô chị lớn ngẩng lên: ``\vie{GATE thì sao?}''

Tôi ngẫm nghĩ. ``\emph{Cổng. Cửa. Lối ra\dots}\ \vie{Ừ, hợp lý.}''

Chúng tôi lắp lần lượt: G - A -T - E.

Ngay khi chữ E cuối cùng gắn vào, mặt sàn bỗng chấn động dữ dội, mạnh hơn bất kỳ rung chuyển nào trước đó. Tôi mất thăng bằng, ngã dúi về phía trước, Kaigou và Suzuki cũng vậy. Những đường nứt chằng chịt chạy khắp sàn như mạng nhện.

*RẮC!*

Một mảng sàn bật tung ra, để lộ một lỗ xoáy màu tím đen đang hút mọi thứ xuống dưới. Nó chính là thứ mà chúng tôi từng thấy khi bước vào The Void.

``\textbf{\vie{Đằng ấy!}}'' tôi hét lên.

Nhưng tôi không cần phải nhắc. Sức hút từ chiếc lỗ ấy đã kéo cả tôi, Kaigou, Suzuki và cả giọng nói đó về phía nó. Chữ cái, bảng lắp, cả không gian quen thuộc dần bị bóp méo, biến dạng như một giấc mơ tan rã.

``\textbf{Chúng ta tìm thấy lối ra rồi!}'' tôi hét, nắm chặt tay Kaigou-chan và kéo cả Suzuki vào lòng, cố giữ ba người lại với nhau.

Nhưng sức hút càng lúc càng mạnh. Mọi thứ tối sầm, ánh sáng biến mất, chỉ còn tiếng hét của chúng tôi vang vọng trong không trung, vang vọng mãi trong vô tận\dots\ rồi lịm hẳn.

*BỤP*

\il

Tôi vẫn còn nhớ rõ cái cảm giác lúc đó, khoảnh khắc cả ba chúng tôi rơi khỏi vùng xoáy ấy, bị nhấn chìm trong màu đen vô tận. Cứ như thể mọi thứ vừa kết thúc, nhưng cũng đồng thời là một khởi đầu.

Phải rồi\dots\ Chúng tôi đã thoát khỏi nơi ``không-có-trò-chơi'' một không gian tưởng chừng vô nghĩa, nhưng lại đầy những câu hỏi, cảm xúc, tiếng cười, và cả một chút\dots\ tai nạn bất ngờ. Nhưng hành trình của chúng tôi vẫn chưa dừng lại ở đó.

Bởi lẽ, khi ánh sáng lại một lần nữa mở ra trước mắt\dots

\dots chúng tôi đã đặt chân đến một thế giới hoàn toàn khác, vừa quen mà vừa lạ. Một nơi đầy sắc màu, đầy giọng nói sôi nổi, và cả những tình huống lố bịch đến mức khó tin.

% Đó là thế giới của NIJISANJI - nơi mà những con người mới xuất hiện, với cá tính kỳ lạ không kém gì những tài năng của tôi. Họ chào đón chúng tôi bằng những biểu cảm rối rắm, tiếng la hét, và sự hỗn loạn không hề xa lạ - giống hệt như mấy trò điên rồ trong \textit{hologra}. Chẳng mấy chốc, chúng tôi đã bị cuốn vào một chuỗi trò trêu chọc, đối đầu, và những trò nghịch phá mang màu sắc rất riêng của nơi này. Một thế giới tưởng chừng như đơn giản, lại khiến tôi bận rộn không kém gì lúc còn ở văn phòng \hll.

% Và sau mười hai ngày kẹt ở đó, là một thế giới hoàn toàn khác, cũng vừa quen vừa lạ.

% Một nơi lạnh lẽo và u ám. Một thị trấn vùng ngoại ô chìm trong khung cảnh hoài cổ của Nhật Bản thời Heisei xen kẽ với thời hiện đại, với những con hẻm tối tăm và ngôi trường Aogami phủ đầy lớp bụi của quá khứ. Nơi đó\dots\ là hololive ERROR.

% Tôi không còn nhớ mình là ai. Kaigou và Suzuki cũng vậy. Chúng tôi, từng người một, phải tìm lại danh tính của chính mình giữa một thực tại đang vỡ vụn, nơi mà một thực thể mang gương mặt Tokino Sora đang dõi theo mọi bước chân của chúng tôi từ trong bóng tối. Và ở trung tâm của vòng lặp vĩnh cửu ấy, chúng tôi đã vô tình hồi sinh Naomi - một cô bé với đôi mắt xanh lục bảo, kẻ đã bị thời gian lãng quên suốt hàng thế kỷ, vẫn đang lặng lẽ đứng dưới tán cây anh đào ngàn năm để chờ đợi một bàn tay chạm tới. Mỗi hành lang, mỗi trang sổ, mỗi bức ảnh\dots\ đều chứa đựng bí mật khiến chúng tôi không thể quay đầu.

% Nhưng đó lại là những câu chuyện mà tôi sẽ kể vào một ngày khác. Sớm thôi.

% Còn bây giờ\dots

% \dots hãy để tôi đưa mọi người quay lại nơi mà tất cả đã kết thúc, nơi mà chúng tôi - Kuon, Kaigou, Suzuki, Game và cả cô bé Naomi - đặt chân đến văn phòng \hll\ vào một ngày thu mờ sương, để tiếp tục một buổi sáng tưởng chừng vô vị\dots

% \dots trước khi cả \hll\ bước ngoặt sang một cuộc hành trình mới.

\end{document}